% Options for packages loaded elsewhere
\PassOptionsToPackage{unicode}{hyperref}
\PassOptionsToPackage{hyphens}{url}
%
\documentclass[
]{article}
\usepackage{amsmath,amssymb}
\usepackage{iftex}
\ifPDFTeX
  \usepackage[T1]{fontenc}
  \usepackage[utf8]{inputenc}
  \usepackage{textcomp} % provide euro and other symbols
\else % if luatex or xetex
  \usepackage{unicode-math} % this also loads fontspec
  \defaultfontfeatures{Scale=MatchLowercase}
  \defaultfontfeatures[\rmfamily]{Ligatures=TeX,Scale=1}
\fi
% \usepackage{lmodern} (commented; using unicode-fix)
% unicode-fix.tex — render raw-Unicode symbols that the Latin Modern *text* font
% lacks, using Computer Modern math/text equivalents. This keeps the established
% Computer Modern look while filling the glyphs xelatex would otherwise drop.
% Included via: pandoc ... --include-in-header=unicode-fix.tex
\usepackage{newunicodechar}
\usepackage{pifont}
% --- Greek letters -> math (Latin Modern Math; consistent with CM body text) ---
\newunicodechar{α}{\ensuremath{\alpha}}
\newunicodechar{β}{\ensuremath{\beta}}
\newunicodechar{γ}{\ensuremath{\gamma}}
\newunicodechar{δ}{\ensuremath{\delta}}
\newunicodechar{ε}{\ensuremath{\varepsilon}}
\newunicodechar{θ}{\ensuremath{\theta}}
\newunicodechar{μ}{\ensuremath{\mu}}
\newunicodechar{ν}{\ensuremath{\nu}}
\newunicodechar{ρ}{\ensuremath{\rho}}
\newunicodechar{σ}{\ensuremath{\sigma}}
\newunicodechar{τ}{\ensuremath{\tau}}
\newunicodechar{φ}{\ensuremath{\varphi}}
\newunicodechar{χ}{\ensuremath{\chi}}
\newunicodechar{ψ}{\ensuremath{\psi}}
\newunicodechar{η}{\ensuremath{\eta}}
\newunicodechar{κ}{\ensuremath{\kappa}}
\newunicodechar{λ}{\ensuremath{\lambda}}
\newunicodechar{π}{\ensuremath{\pi}}
\newunicodechar{ω}{\ensuremath{\omega}}
\newunicodechar{Ḣ}{\ensuremath{\dot H}}
\newunicodechar{≪}{\ensuremath{\ll}}
\newunicodechar{≫}{\ensuremath{\gg}}
\newunicodechar{⟺}{\ensuremath{\Longleftrightarrow}}
\newunicodechar{⟹}{\ensuremath{\Longrightarrow}}
\newunicodechar{⇔}{\ensuremath{\Leftrightarrow}}
\newunicodechar{⟨}{\ensuremath{\langle}}
\newunicodechar{⟩}{\ensuremath{\rangle}}
\newunicodechar{⁰}{\textsuperscript{0}}
\newunicodechar{¹}{\textsuperscript{1}}
\newunicodechar{²}{\textsuperscript{2}}
\newunicodechar{³}{\textsuperscript{3}}
\newunicodechar{⁴}{\textsuperscript{4}}
\newunicodechar{⁵}{\textsuperscript{5}}
\newunicodechar{⁶}{\textsuperscript{6}}
\newunicodechar{⁷}{\textsuperscript{7}}
\newunicodechar{⁸}{\textsuperscript{8}}
\newunicodechar{⁹}{\textsuperscript{9}}
% --- Relations / operators -> math ---
\newunicodechar{↔}{\ensuremath{\leftrightarrow}}
\newunicodechar{⇐}{\ensuremath{\Leftarrow}}
\newunicodechar{⇒}{\ensuremath{\Rightarrow}}
\newunicodechar{∂}{\ensuremath{\partial}}
\newunicodechar{∅}{\ensuremath{\varnothing}}
\newunicodechar{∈}{\ensuremath{\in}}
\newunicodechar{∼}{\ensuremath{\sim}}
\newunicodechar{≈}{\ensuremath{\approx}}
\newunicodechar{≠}{\ensuremath{\neq}}
\newunicodechar{≤}{\ensuremath{\leq}}
\newunicodechar{≥}{\ensuremath{\geq}}
\newunicodechar{⊂}{\ensuremath{\subset}}
\newunicodechar{⊃}{\ensuremath{\supset}}
% --- Special ---
\newunicodechar{ℏ}{\ensuremath{\hbar}}
\newunicodechar{ℤ}{\ensuremath{\mathbb{Z}}}
% --- Superscript modifier letters (footnote-style markers) -> text superscript ---
\newunicodechar{ʰ}{\textsuperscript{h}}
\newunicodechar{ʲ}{\textsuperscript{j}}
\newunicodechar{ᵃ}{\textsuperscript{a}}
\newunicodechar{ᵇ}{\textsuperscript{b}}
\newunicodechar{ᵈ}{\textsuperscript{d}}
\newunicodechar{ᵉ}{\textsuperscript{e}}
\newunicodechar{ᵍ}{\textsuperscript{g}}
\newunicodechar{ᶠ}{\textsuperscript{f}}
\newunicodechar{⁻}{\textsuperscript{\ensuremath{-}}}
% --- Subscript digits -> text subscript ---
\newunicodechar{₀}{\textsubscript{0}}
\newunicodechar{₁}{\textsubscript{1}}
\newunicodechar{₂}{\textsubscript{2}}
\newunicodechar{₃}{\textsubscript{3}}
\newunicodechar{₄}{\textsubscript{4}}
\newunicodechar{₈}{\textsubscript{8}}
% --- Checkmarks ---
\newunicodechar{✓}{\ding{51}}
\newunicodechar{✗}{\ding{55}}

\ifPDFTeX\else
  % xetex/luatex font selection
\fi
% Use upquote if available, for straight quotes in verbatim environments
\IfFileExists{upquote.sty}{\usepackage{upquote}}{}
\IfFileExists{microtype.sty}{% use microtype if available
  \usepackage[]{microtype}
  \UseMicrotypeSet[protrusion]{basicmath} % disable protrusion for tt fonts
}{}
\makeatletter
\@ifundefined{KOMAClassName}{% if non-KOMA class
  \IfFileExists{parskip.sty}{%
    \usepackage{parskip}
  }{% else
    \setlength{\parindent}{0pt}
    \setlength{\parskip}{6pt plus 2pt minus 1pt}}
}{% if KOMA class
  \KOMAoptions{parskip=half}}
\makeatother
\usepackage{xcolor}
\usepackage{longtable,booktabs,array}
\usepackage{calc} % for calculating minipage widths
% Correct order of tables after \paragraph or \subparagraph
\usepackage{etoolbox}
\makeatletter
\patchcmd\longtable{\par}{\if@noskipsec\mbox{}\fi\par}{}{}
\makeatother
% Allow footnotes in longtable head/foot
\IfFileExists{footnotehyper.sty}{\usepackage{footnotehyper}}{\usepackage{footnote}}
\makesavenoteenv{longtable}
\setlength{\emergencystretch}{3em} % prevent overfull lines
\providecommand{\tightlist}{%
  \setlength{\itemsep}{0pt}\setlength{\parskip}{0pt}}
\setcounter{secnumdepth}{-\maxdimen} % remove section numbering
\ifLuaTeX
  \usepackage{selnolig}  % disable illegal ligatures
\fi
\IfFileExists{bookmark.sty}{\usepackage{bookmark}}{\usepackage{hyperref}}
\IfFileExists{xurl.sty}{\usepackage{xurl}}{} % add URL line breaks if available
\urlstyle{same}
\hypersetup{
  hidelinks,
  pdfcreator={LaTeX via pandoc}}

\author{}
\date{}

\begin{document}

\hypertarget{non-markovian-dynamics-in-computational-biology}{%
\section{Non-Markovian Dynamics in Computational
Biology}\label{non-markovian-dynamics-in-computational-biology}}

\hypertarget{hidden-sector-memory-and-the-methodological-limits-of-single-cell-analysis}{%
\subsubsection{Hidden-Sector Memory and the Methodological Limits of
Single-Cell
Analysis}\label{hidden-sector-memory-and-the-methodological-limits-of-single-cell-analysis}}

\textbf{Author:} Alex Maybaum \textbf{Date:} May 2026
\textbf{Classification:} Computational Biology / Bioinformatics

\begin{center}\rule{0.5\linewidth}{0.5pt}\end{center}

\hypertarget{abstract}{%
\subsection{Abstract}\label{abstract}}

The single-cell bioinformatics literature has documented a series of
methodological puzzles over the past five years: RNA velocity methods
that reverse trajectories on validated developmental data; gene
regulatory network inference accuracy that ``marginally exceeds random
predictions'' despite extensive method development; deep-learning
foundation models (scGPT, Geneformer, scFoundation) that systematically
underperform simple linear baselines on perturbation prediction;
multi-omics integration methods that produce inconsistent results across
modalities. Each puzzle has been treated as a methodological problem to
be solved with better algorithms, more data, or different model
architectures.

This paper applies the Observational Incompleteness (OI) framework
{[}1{]} to these puzzles. The framework's C1--C3 conditions ---
non-trivial coupling (C1), slow-bath memory (C2), and high-capacity
hidden sector (C3) --- produce P-indivisible (non-Markovian) dynamics
whenever a fast visible sector is coupled to a slow high-capacity hidden
sector. The companion Medicine paper {[}2{]} establishes that chromatin
layers satisfy C1--C3 in biological systems, with DNA methylation,
histone modifications, and chromatin state functioning as the slow
hidden sector with timescales spanning seconds (chromatin binding) to
generations (DNA methylation maintenance). This paper shows that the
same structural conditions explain the bioinformatics field's documented
failure modes: each puzzle has the same underlying cause --- methods
assume Markovian dynamics on systems with substantial hidden-sector
memory, and this assumption violation produces characteristic systematic
errors.

We develop the analysis across five domains where the framework's
predictions match documented empirical failure modes: trajectory
inference (RNA velocity assumptions vs.~multi-omic velocity
improvements); gene regulatory network inference (correlation-based
methods outperforming causal methods, multi-omic methods outperforming
RNA-only); perturbation prediction (deep-learning foundation models
underperforming linear baselines, knowledge-graph methods improving on
data-driven methods); multi-omics integration (static factor analysis
vs.~emerging timescale-aware methods); and cancer methylome
classification (95\% subclass accuracy of methylation-based tumor
classifiers reflecting substratum-preservation in cancer). We also
extend the framework's reader-writer disorder class {[}Medicine §10.7{]}
into bioinformatic predictions, including a testable claim about variant
pathogenicity stratification for reader/writer/eraser genes.

The structural content of the framework's application to computational
biology is informational rather than algorithmic. Standard methodology
assumes that more data and more complex models will overcome any
prediction task; the framework identifies an information-theoretic
ceiling that no transcriptome-only model can exceed, no matter how
scaled. The methodological roadmap forward is to include hidden-sector
information explicitly --- through multi-modal data (chromatin
accessibility, histone marks, methylation), knowledge graphs (encoding
hidden regulatory structure), or explicit memory state representations
--- rather than to scale transcriptome-only architectures further.

\begin{center}\rule{0.5\linewidth}{0.5pt}\end{center}

\hypertarget{introduction}{%
\subsection{1. Introduction}\label{introduction}}

\hypertarget{the-bioinformatics-fields-persistent-puzzles}{%
\subsubsection{1.1 The bioinformatics field's persistent
puzzles}\label{the-bioinformatics-fields-persistent-puzzles}}

Single-cell transcriptomics has revolutionized biology over the past
decade, producing datasets of millions of cells with detailed gene
expression profiles. The computational methodology has evolved rapidly:
dimensionality reduction (UMAP, t-SNE), clustering (Leiden, Louvain),
trajectory inference (Monocle, Slingshot, scVelo), gene regulatory
network inference (SCENIC, GENIE3), perturbation modeling (CPA, GEARS,
scGPT), and multi-omics integration (MOFA, Seurat, LIGER). Each new
method generation has promised improved accuracy and broader
applicability, and the field's expected trajectory has been continuous
improvement.

A different pattern has emerged. Recent benchmarking studies have
documented systematic failures across multiple method classes:

\begin{itemize}
\item
  \textbf{RNA velocity} methods (velocyto, scVelo) ``fail to identify,
  and even reverses the trajectory'' on validated developmental datasets
  {[}3{]}. The failures are reproducible across at least six independent
  datasets {[}3{]}.
\item
  \textbf{Gene regulatory network inference} accuracy ``marginally
  exceeds random predictions'' despite a decade of method development
  {[}4{]}. Causal methods consistently fail to outperform
  correlation-based baselines --- a ``persistent puzzle which casts
  doubt on the value of causality for this task'' {[}5{]}.
\item
  \textbf{Deep-learning foundation models} for single-cell biology
  (scGPT, Geneformer, scFoundation, scPRINT) ``none outperformed
  {[}linear{]} baselines'' for predicting transcriptome changes after
  perturbations {[}6{]}. Zero-shot evaluations show foundation models
  ``perform worse than selecting highly variable genes and using more
  established methods such as Harmony and scVI in cell type clustering''
  {[}7{]}.
\item
  \textbf{Multi-omics integration} methods (MOFA, MOFA+) treat different
  molecular modalities as parallel measurements of the same cell state,
  despite the fact that ``the timescale of (phospho)proteome changes
  falls into a minute regime, whereas most significant transcriptome
  changes are observed after hours'' {[}8{]}.
\end{itemize}

Each puzzle has been addressed within the field as a methodological
problem to be solved with better algorithms. The proposed solutions
follow a common pattern: incorporate additional data modalities
(chromatin accessibility, ATAC-seq, methylation), incorporate prior
knowledge (knowledge graphs, motif libraries), or move to more complex
architectures (transformers, foundation models, graph neural networks).
Some of these solutions help; others do not. The field lacks a
structural framework for predicting \emph{which} methodological
improvements will help and \emph{why}.

\hypertarget{the-oi-framework-and-its-c1c3-conditions}{%
\subsubsection{1.2 The OI framework and its C1--C3
conditions}\label{the-oi-framework-and-its-c1c3-conditions}}

The Observational Incompleteness (OI) framework {[}1{]} establishes a
characterization theorem for embedded observation. The theorem states:
any embedded observer whose dynamics satisfies three structural
conditions --- C1 (non-trivial coupling between visible and hidden
sectors), C2 (slow-bath memory, \(\tau_B \gg \tau_S\)), and C3
(high-capacity hidden sector) --- necessarily experiences P-indivisible
dynamics: history-dependent transition probabilities arising from
information stored in the hidden sector and returned on subsequent
interactions. The theorem is scale-independent: it applies to any system
with the appropriate architecture, from quantum mechanics emerging at
the cosmological horizon to biological signaling cascades.

The companion Medicine paper {[}2{]} establishes that biological memory
systems satisfy C1--C3 at multiple molecular layers. The chromatin layer
is the most directly relevant for computational biology: DNA methylation
patterns, histone modifications, and chromatin-state assignments form a
slow hidden sector with substantial \(\tau_B\) relative to
transcriptional dynamics, and the chromatin state space is exponentially
large (over \(4^{N_{\text{CpG}}}\) for \(N_{\text{CpG}} \sim 28\)
million CpG sites in the human genome). Cellular dynamics --- the
trajectory a cell takes through state space during differentiation,
perturbation response, or homeostasis --- depends on both visible-sector
dynamics (transcription, signaling) and hidden-sector state (chromatin
memory).

The framework's bioinformatic prediction follows from the theorem:
methods that infer cellular dynamics from visible-sector data alone
(transcriptome) must systematically miss the contribution of
hidden-sector memory. The errors are not random noise; they have a
specific structure: directional biases, reduced accuracy in transitions
where chromatin state changes, and information-theoretic ceilings on
what visible-sector data can recover. Methods that include hidden-sector
information explicitly (multi-modal data, knowledge graphs, prior
biological information) should systematically outperform pure
visible-sector methods on these tasks.

\hypertarget{thesis-and-structure}{%
\subsubsection{1.3 Thesis and structure}\label{thesis-and-structure}}

This paper develops the framework's prediction across five active
methodological domains in single-cell bioinformatics. Each section
follows a common structure: (i) the documented field-level puzzle or
failure mode; (ii) the framework's structural explanation in terms of
C1--C3 dynamics; (iii) the empirically observed methodological
improvements that align with the framework's prediction.

§2 addresses trajectory inference, from RNA velocity's documented
failures to multi-omic velocity's improvements. §3 addresses gene
regulatory network inference, including the field's ``persistent
puzzle'' of causal methods failing to outperform correlation. §4
addresses perturbation prediction, where foundation models
systematically underperform simple baselines. §5 addresses multi-omics
integration, where static methods struggle with the timescale hierarchy
across modalities. §6 addresses cancer methylome classification, where
the framework's substratum-preservation prediction matches the entire
field's diagnostic paradigm. §7 develops a new testable prediction about
variant pathogenicity stratification for reader/writer/eraser genes. §8
extends the framework's reader-writer disorder class {[}Medicine
§10.7{]} into bioinformatic patterns. §9 synthesizes the cross-domain
findings. §10 catalogs open questions and limits.

The paper's contribution is not new algorithms. It is a structural
framework for understanding why current methods fail in specific
patterns, and for predicting which methodological improvements will
work. The information-theoretic ceiling identified here has practical
consequences: it explains why scaling foundation models has not produced
the expected improvements, and predicts that future progress requires
architectural changes that incorporate hidden-sector information rather
than scaling visible-sector models further.

\begin{center}\rule{0.5\linewidth}{0.5pt}\end{center}

\hypertarget{trajectory-inference}{%
\subsection{2. Trajectory Inference}\label{trajectory-inference}}

\hypertarget{the-rna-velocity-paradigm-and-its-assumptions}{%
\subsubsection{2.1 The RNA velocity paradigm and its
assumptions}\label{the-rna-velocity-paradigm-and-its-assumptions}}

RNA velocity, introduced by La Manno et al.~(2018) {[}9{]} and
operationalized in the velocyto and scVelo tools, infers the direction
of cellular state change from the ratio of unspliced to spliced mRNA in
single-cell RNA sequencing data. The method's appeal is that it extracts
dynamic information from static snapshots: by measuring the relative
abundance of pre-mRNA (unspliced, recently transcribed) and mature mRNA
(spliced, accumulated), it estimates whether a gene is in an induction
phase (more unspliced relative to expected steady state) or repression
phase (more spliced relative to expected steady state). Aggregating
across genes produces a velocity vector for each cell, which can be
visualized as a flow field over the dimensionally-reduced state space.

The original velocyto method {[}9{]} assumes (i) constant rates of
transcription and degradation, (ii) a global splicing rate shared across
genes, (iii) the dynamics has reached steady-state equilibrium in the
induction phase, and (iv) gene-wise independence. The scVelo extension
{[}10{]} relaxes (i) and (iii) but retains the assumption that
within-gene dynamics is governed by a closed kinetic equation depending
only on current mRNA levels. Both methods treat the gene-state dynamics
as Markovian: the probability of state change at time \(t + \Delta t\)
depends only on the state at time \(t\), not on the cell's history.

\hypertarget{documented-failure-modes}{%
\subsubsection{2.2 Documented failure
modes}\label{documented-failure-modes}}

Gorin, Pachter, and Yosef (2022) {[}3{]} published a detailed analysis
of RNA velocity methods in PLOS Computational Biology, demonstrating
that scVelo ``fails to identify, and even reverses the trajectory'' in
multiple validated developmental datasets. The reversed trajectories
produce qualitatively wrong biological conclusions: cell types
identified as progenitors are actually terminal states, and the inferred
differentiation direction is opposite to the ground truth. The failures
are documented across at least six independent datasets {[}3, 11--14{]}.

Barile et al.~(2021) {[}11{]} and subsequent analyses identify the
proximate cause: when cell populations have heterogeneous kinetics
(different genes operating at different rates) or when the experiment
captures cells not in steady state, scVelo's parameter estimates are
biased in a way that systematically reverses the inferred direction. The
DeepVelo paper {[}12{]} notes that ``these assumptions are often
violated in complex biological systems and bring about limitations in
downstream applications, particularly when cell states are partially
observed or undergo transcription dynamics more complex than the
steady-state pattern.''

The pattern is consistent: RNA velocity methods work for simple linear
trajectories with homogeneous kinetics and fail for complex trajectories
with heterogeneous dynamics. The methodological response has been to
develop progressively more flexible parameter estimation (LatentVelo
{[}15{]}, veloVI {[}16{]}, DeepVelo {[}12{]}). These improvements help
marginally but do not eliminate the failure mode: the underlying issue
is structural, not parametric.

\hypertarget{the-frameworks-structural-explanation}{%
\subsubsection{2.3 The framework's structural
explanation}\label{the-frameworks-structural-explanation}}

The OI framework predicts the RNA velocity failure mode structurally.
Gene expression in single cells is governed by the joint dynamics of the
transcriptional machinery (visible sector) and the chromatin state
(hidden sector). Chromatin state determines which genes are accessible
for transcription, at what rates, and with what kinetic parameters;
chromatin state itself evolves on slower timescales than transcription,
with histone modifications turning over on minutes-to-hours scales and
DNA methylation on days-to-weeks scales {[}Medicine §9.1{]}. The full
dynamics is non-Markovian in the gene-expression representation: knowing
the current spliced/unspliced ratios is insufficient to predict the next
state, because chromatin memory contributes information not captured by
mRNA levels alone.

The C1--C3 architecture applies directly: gene transcription (fast,
visible) is coupled (C1) to chromatin state (slow, C2) with
exponentially large chromatin state space (C3). The framework's
characterization theorem {[}1{]} then implies that RNA velocity inferred
from gene-expression data alone is necessarily P-indivisible: the
marginal dynamics on the gene-expression subspace cannot be described by
a Markovian kinetic equation, and methods assuming such description will
produce biased estimates.

The bias is not arbitrary. Its sign depends on the relationship between
chromatin state and transcriptional kinetics. When chromatin transitions
are slow relative to the measurement timescale (typical for
differentiation), the chromatin state is approximately frozen, and
gene-expression dynamics within a chromatin state can be approximately
Markovian --- RNA velocity works. When chromatin transitions are
comparable to or faster than the measurement timescale (typical for
cell-fate transitions, perturbation responses, and complex
differentiation), the chromatin state changes substantially within the
observation window, and gene-expression dynamics is non-Markovian ---
RNA velocity fails, with specific systematic biases depending on the
chromatin transition pattern.

This matches the empirical pattern in Gorin et al.~and the surrounding
literature: simple linear trajectories with homogeneous kinetics
succeed; complex trajectories with heterogeneous dynamics fail. The
framework explains why this is the structural pattern rather than a
methodological accident.

\hypertarget{multi-omic-velocity-as-empirical-convergence}{%
\subsubsection{2.4 Multi-omic velocity as empirical
convergence}\label{multi-omic-velocity-as-empirical-convergence}}

The methodological response that does work is to include chromatin
accessibility data explicitly. MultiVelo {[}17{]}, introduced by Li et
al.~(2022) in Nature Biotechnology, ``extends the RNA velocity framework
to incorporate epigenomic data'' and ``improves the accuracy of cell
fate prediction compared to velocity estimates from RNA only.'' The
model includes both transcriptional dynamics (spliced/unspliced ratios)
and chromatin accessibility dynamics, fitting a coupled differential
equation system. By estimating switch times for chromatin
opening/closing alongside transcription rates, MultiVelo recovers two
distinct classes of genes (chromatin closes before vs.~after
transcription ceases) and four types of cell states distinguishing
coupled and decoupled chromatin-transcription dynamics.

The methodological trajectory has continued. mmVelo {[}18{]} (2024)
extends multi-modal velocity estimation to single-peak resolution rather
than gene-aggregated chromatin signals. MoFlow {[}19{]} (Nature
Communications 2025) integrates multi-omic data for ``locally adaptive,
single-cell-resolved velocities that illuminate fine-scale regulatory
dynamics.'' Each successive method incorporates more hidden-sector
information (chromatin accessibility at higher resolution, additional
modalities, locally adaptive dynamics) and improves performance.

This trajectory matches the framework's prediction. The methods that
work are the methods that include the hidden sector explicitly; the
methods that fail are the methods that assume Markovian dynamics on the
visible sector alone. The empirical convergence over the 2022--2025
period --- from RNA-only methods through gene-level multi-omic methods
to single-peak resolution multi-omic methods --- represents the field's
gradual adoption of the framework's structural prediction without
explicit reference to the underlying theorem.

\hypertarget{practical-consequences}{%
\subsubsection{2.5 Practical
consequences}\label{practical-consequences}}

The framework provides three practical predictions for trajectory
inference methodology:

(P1) RNA-only trajectory methods will continue to fail for systems with
substantial chromatin dynamics during the observation window. Parameter
refinements within the RNA-only paradigm (improved kinetic models,
deeper neural networks on transcriptome alone) cannot fix this; the
limit is information-theoretic.

(P2) Multi-modal methods that include chromatin accessibility will
outperform RNA-only methods on differentiation, perturbation response,
and cell-fate transition tasks. The improvement should scale with the
relative timescale of chromatin transitions to transcription: faster
chromatin dynamics relative to transcription means larger improvement
from including chromatin data.

(P3) The optimal data modality combination depends on the timescale of
the biological process being studied. For processes with \(\tau_B\) at
the histone-modification timescale (minutes-to-hours), histone-mark
profiling (CUT\&Tag, CUT\&RUN) is more informative than DNA methylation.
For processes at the methylation timescale (days-to-weeks), methylation
profiling is more informative. For long-term memory (developmental
trajectories spanning generations), DNA methylation and chromatin
architecture (Hi-C) become the dominant relevant modalities. Mixing
modalities should be calibrated to the biological process's
characteristic \(\tau_B\) rather than included indiscriminately.

The third prediction is testable through benchmarking: a study comparing
multi-modal velocity methods across biological processes with different
characteristic timescales should show the relative performance of each
modality combination tracking the framework's prediction.

\begin{center}\rule{0.5\linewidth}{0.5pt}\end{center}

\hypertarget{gene-regulatory-network-inference}{%
\subsection{3. Gene Regulatory Network
Inference}\label{gene-regulatory-network-inference}}

\hypertarget{the-grn-inference-paradigm-and-the-persistent-puzzle}{%
\subsubsection{3.1 The GRN inference paradigm and the ``persistent
puzzle''}\label{the-grn-inference-paradigm-and-the-persistent-puzzle}}

Gene regulatory networks (GRNs) describe the directed influence of
transcription factors on target genes. Inferring GRNs from single-cell
expression data is a foundational task in computational biology, with
applications ranging from drug-target identification to mechanistic
understanding of differentiation programs. The methodology has produced
a series of method generations: correlation-based methods (Pearson,
mutual information), ensemble methods (GENIE3, GRNBoost2),
expression-cell-type specific methods (SCENIC), causal methods (PC
algorithm, NOTEARS, DeepSEM), and most recently deep-learning foundation
models (scGPT, Geneformer) repurposed for GRN inference.

Despite the methodological investment, the empirical state of GRN
inference accuracy is sobering. Yuan et al.~(Nature Biotechnology 2024)
{[}20{]} state: ``Despite extensive efforts, GRN inference accuracy has
remained disappointingly low, marginally exceeding random predictions.''
Passemiers et al.~(2025) {[}5{]} document a ``persistent puzzle which
casts doubt on the value of causality for this task'': causal methods
consistently fail to outperform correlation-based baselines in realistic
benchmarks, despite their theoretical advantages.

Multiple comprehensive benchmarking studies have confirmed this pattern.
McCalla et al.~{[}21{]} evaluated 13 GRN inference methods on 7
published datasets using diverse gold standards; the methods rank
similarly across gold standards, with no method showing decisive
advantage. Nguyen et al.~{[}22{]} benchmarked 15 methods on 139
simulated datasets and found extensive failure modes related to dropout
rate and sparsity. The conclusion across the literature is consistent:
GRN inference from single-cell data is hard, and current methods perform
near baseline.

\hypertarget{why-correlation-outperforms-causation-the-frameworks-prediction}{%
\subsubsection{3.2 Why correlation outperforms causation: the
framework's
prediction}\label{why-correlation-outperforms-causation-the-frameworks-prediction}}

The persistent puzzle has a natural framework explanation. Causal
methods assume that the joint distribution over gene expression captures
the system's dynamics: from sufficient samples, the conditional
independence structure reveals the causal graph. This assumption is
valid for systems where the visible-sector dynamics is autonomous ---
where current gene expression states fully determine future states up to
noise.

The framework predicts this assumption fails for cellular systems. Gene
regulation depends on chromatin state, transcription factor protein
levels, post-translational modifications, and other variables not
captured in scRNA-seq data. The marginal distribution over gene
expression is the result of integrating out these hidden variables, and
the conditional independence structure of the marginal distribution does
\emph{not} match the causal graph of the underlying dynamics. Causal
inference methods recover the marginal structure, not the underlying
causal structure, because the data does not contain enough information
to distinguish them.

Correlation-based methods make a weaker claim: they identify pairs of
genes whose expression covaries, without claiming causal direction. The
covariance pattern is robust to hidden-variable integration in a way
that the causal structure is not. Correlation-based methods perform
competitively because they target a tractable quantity (the marginal
covariance) rather than an intractable one (the marginal causal
structure).

This is the framework's prediction: in non-Markovian systems where the
visible sector does not contain the full dynamical information,
correlation-based methods will systematically outperform causal methods
that assume Markovian dynamics on the visible sector. The ``persistent
puzzle'' is the empirical signature of this structural prediction.

\hypertarget{multi-modal-methods-as-the-predicted-improvement}{%
\subsubsection{3.3 Multi-modal methods as the predicted
improvement}\label{multi-modal-methods-as-the-predicted-improvement}}

The methodological response that improves GRN inference accuracy is to
incorporate additional data modalities. SCENIC+ {[}23{]}, introduced by
González-Blas et al.~(Nature Methods 2023), extends the original SCENIC
pipeline by incorporating chromatin accessibility data (scATAC-seq)
alongside gene expression. The method identifies enhancer-driven gene
regulatory networks by predicting genomic enhancers, linking them to
candidate transcription factors via motif analysis, and linking
enhancers to target genes via expression correlation across cells.
SCENIC+ improves substantially over RNA-only SCENIC and over earlier
methods that used only expression data.

LINGER {[}24{]}, introduced by Yuan et al.~(Nature Biotechnology 2024),
takes the multi-modal approach further by incorporating ``atlas-scale
external bulk data across diverse cellular contexts and prior knowledge
of transcription factor motifs as a manifold regularization.'' The
method combines paired scRNA-seq + scATAC-seq data with external bulk
reference data, encoding prior biological knowledge about transcription
factor binding preferences. The performance improvement is substantial:
LINGER's GRN predictions outperform SCENIC+ and earlier methods across
multiple benchmarks.

The trajectory matches the framework's prediction. RNA-only methods are
near random baseline. Adding chromatin accessibility (substratum-level
visible) helps. Adding prior biological knowledge (encoded hidden
structure) helps further. Each addition incorporates more hidden-sector
information, and each addition improves accuracy. The framework's
structural prediction --- that GRN inference accuracy scales with how
much hidden-sector information is included in the data --- is
empirically realized in the methodological development trajectory of the
past five years.

\hypertarget{information-theoretic-ceiling}{%
\subsubsection{3.4 Information-theoretic
ceiling}\label{information-theoretic-ceiling}}

The framework predicts more than just a relative improvement from
multi-modal data. It predicts an information-theoretic ceiling: no
method using only scRNA-seq data can recover the true causal GRN,
regardless of model complexity or training data scale. This is the
structural content of P-indivisibility: information that has been
integrated out cannot be recovered from the marginal observations.

The recent attempts to apply deep-learning foundation models to GRN
inference test this prediction directly. scGPT {[}25{]} and Geneformer
{[}26{]}, trained on tens of millions of single cells, attempt to learn
gene-gene relationships from transcriptome data alone. Multiple
independent evaluations have shown that these models do not improve
substantially over simpler correlation-based methods on GRN inference
(§4 develops this for perturbation prediction; similar results hold for
GRN tasks). The foundation models' attention patterns ``capture
co-expression rather than unique regulatory signal'' {[}27{]} --- they
recover the marginal correlation structure that simpler methods also
recover, without accessing the underlying regulatory causal structure.

This is the information-theoretic ceiling in action. Scaling
transcriptome-only foundation models does not break through the ceiling;
the ceiling is structural, not computational. The methodological roadmap
forward requires architectural changes that incorporate hidden-sector
information, not further scaling of visible-sector models.

\hypertarget{practical-consequences-1}{%
\subsubsection{3.5 Practical
consequences}\label{practical-consequences-1}}

Three framework-grounded predictions for GRN inference methodology:

(P4) GRN inference from scRNA-seq alone is near its
information-theoretic ceiling. Further methodological development that
stays within RNA-only data will produce marginal improvements, not the
order-of-magnitude gains the field has been seeking.

(P5) Multi-modal methods (RNA + ATAC, RNA + methylation, RNA + Hi-C)
will systematically outperform RNA-only methods. The relative gain
depends on which hidden-sector modality is most informative for the
specific regulatory structure being inferred.

(P6) Foundation models applied to scRNA-seq alone will not exceed the
information-theoretic ceiling, regardless of scale. Foundation models
trained on multi-modal data (paired RNA + ATAC, multi-omics atlases)
should outperform single-modality foundation models on GRN inference
tasks. This is testable: compare scGPT trained on RNA-only versus a
foundation model of comparable scale trained on multi-modal data.

The third prediction is testable in the near term as multi-modal
foundation models become available.

\begin{center}\rule{0.5\linewidth}{0.5pt}\end{center}

\hypertarget{perturbation-prediction}{%
\subsection{4. Perturbation Prediction}\label{perturbation-prediction}}

\hypertarget{the-perturbation-prediction-paradigm}{%
\subsubsection{4.1 The perturbation prediction
paradigm}\label{the-perturbation-prediction-paradigm}}

Single-cell perturbation experiments --- CRISPR knockouts (Perturb-seq,
CROP-seq), CRISPRa/i activation/repression (Replogle et al.~2022
{[}28{]}), drug treatments (sci-Plex), and combinatorial perturbations
(Norman et al.~2019 {[}29{]}) --- produce datasets where individual
cells are labeled by the perturbation they received and the resulting
transcriptional state is measured. The computational task is to predict
the transcriptional state for \emph{unseen} perturbations: combinations
of perturbations not in the training data, perturbations applied to new
cell types, or perturbations with new drugs.

This task has attracted substantial methodological investment.
Generative models (CPA {[}30{]}, chemCPA {[}31{]}) learn compositional
representations of perturbation effects. Graph neural networks (GEARS
{[}32{]}) incorporate Gene Ontology and protein-protein interaction
graphs as priors. Optimal transport methods (CellOT {[}33{]}) match
unperturbed and perturbed distributions. Most recently, single-cell
foundation models (scGPT {[}25{]}, Geneformer {[}26{]}, scFoundation
{[}34{]}, scPRINT) have been adapted for perturbation prediction tasks,
with the expectation that pretraining on tens of millions of cells would
enable strong zero-shot and few-shot performance.

\hypertarget{the-empirical-failure-of-deep-learning-on-perturbation-prediction}{%
\subsubsection{4.2 The empirical failure of deep learning on
perturbation
prediction}\label{the-empirical-failure-of-deep-learning-on-perturbation-prediction}}

Ahlmann-Eltze et al.~(Nature Methods 2025) {[}6{]} published a direct
evaluation of this expectation. The study compared five foundation
models --- scGPT, Geneformer, scFoundation, scPRINT, and Universal Cell
Embeddings --- and two additional deep-learning models against
deliberately simple linear baselines for predicting transcriptome
changes after single or double perturbations. The result: ``None
outperformed the baselines.'' The simplest baseline --- predicting the
mean response across the training set --- was competitive with or
superior to every deep-learning method evaluated.

The result has been replicated. The geneRNIB benchmark (Passemiers et
al.~2025) {[}5{]} confirms that deep-learning methods do not
consistently outperform linear baselines on perturbation prediction
tasks. Independent analyses of scGPT's attention mechanisms {[}27{]}
show that the attention patterns ``capture co-expression rather than
unique regulatory signal'' --- the model has learned the gene-gene
correlation structure that simpler methods also capture, without
accessing additional information.

The pattern is structural, not technical. The foundation models are
well-implemented, well-trained, and demonstrably effective on tasks
where the visible-sector data contains sufficient information (cell type
annotation with appropriate fine-tuning, batch correction in some
settings). The failure is specific to tasks where the missing
information matters: predicting how cells will respond to perturbations
not seen in training requires knowing what the cell can become, which
depends on chromatin state, transcription factor activity, and other
variables not in the transcriptome.

\hypertarget{the-frameworks-structural-prediction}{%
\subsubsection{4.3 The framework's structural
prediction}\label{the-frameworks-structural-prediction}}

The OI framework's information-theoretic constraint produces a hard
upper bound on what transcriptome-only perturbation models can achieve.
The argument is direct: the cell's response to a perturbation depends on
(i) the current transcriptional state, (ii) the chromatin state
determining which genes are accessible for activation, (iii) the protein
levels of transcription factors not visible in mRNA data, and (iv) the
metabolic state. Of these, (i) is in the data and (ii)--(iv) are not.
The function mapping perturbation × cell state to response state is
therefore a function of more variables than the data contains; the
mapping cannot be learned from data alone, no matter how much data is
available, because the inverse problem is structurally under-determined.

The constraint is information-theoretic in the strict sense: the mutual
information between the visible-sector data and the perturbation
response is bounded above by a quantity that does not scale with dataset
size. Scaling the model and the data improves estimation of the
conditional mean response (which is why deep-learning models do achieve
some performance), but cannot improve prediction of the response's
\emph{variance}, \emph{outlier distribution}, or \emph{cell-specific
deviations} --- all of which depend on hidden-sector variables.

The Ahlmann-Eltze result has a specific structural interpretation in
this framework: when the baseline (predict the mean) is competitive with
foundation models, it indicates that the foundation models are also
predicting close to the conditional mean response, because the
conditional mean is the limit of what transcriptome data supports. Deep
learning architectures do not extract additional information from the
visible sector because there is no additional information to extract.

\hypertarget{knowledge-graph-methods-as-the-predicted-improvement}{%
\subsubsection{4.4 Knowledge-graph methods as the predicted
improvement}\label{knowledge-graph-methods-as-the-predicted-improvement}}

The methods that do exceed simple baselines incorporate explicit
hidden-sector information through external knowledge. GEARS {[}32{]}
uses Gene Ontology graphs and protein-protein interaction networks as
structural priors during training; the prior knowledge constrains the
prediction in ways that pure data-driven learning cannot. PerturbNet
{[}35{]} (Mol Syst Biol 2025) achieves the highest R² values across
multiple benchmarks by separating the drug representation network
(trained on chemical structure) from the cell representation network
(trained on expression), allowing the model to compose information from
disjoint training sources. Both methods incorporate information that is
not in the scRNA-seq data --- chemical structure, ontological gene
relationships, prior regulatory knowledge --- and both outperform pure
data-driven foundation models.

The methodological trajectory is consistent with the framework: the
methods that work are the methods that include hidden-sector information
explicitly through external knowledge sources, and the methods that fail
are the methods that try to extract that information from transcriptome
data alone. The ``deep learning fails to beat baselines'' puzzle is the
framework's information-theoretic ceiling made empirically manifest.

\hypertarget{compositional-generalization-and-what-scaling-cannot-fix}{%
\subsubsection{4.5 Compositional generalization and what scaling cannot
fix}\label{compositional-generalization-and-what-scaling-cannot-fix}}

A subtler prediction from the framework concerns compositional
generalization --- the ability to predict combinations of perturbations
from training data containing only single perturbations. Standard
intuition is that with enough data and a sufficiently flexible model,
compositional structure should be learnable. The framework's prediction
is more restrictive: compositional generalization works only when the
perturbations act on visible-sector variables independently, and fails
when perturbations interact through hidden-sector variables (chromatin
state, transcription factor competition, etc.).

This is testable. Norman et al.~(2019) {[}29{]} published a
combinatorial perturbation dataset with 287 single-gene knockdowns and
131 pairs of knockdowns. The framework predicts that pairs where both
genes act on the same chromatin domain (likely to interact through
hidden-sector variables) should be harder to predict from single-gene
training data than pairs where the genes act on disjoint pathways. The
empirical pattern in compositional generalization benchmarks should
track the chromatin co-localization of the perturbed genes, not just
their expression patterns. This prediction has not been directly tested
in the existing benchmarking literature.

\hypertarget{practical-consequences-2}{%
\subsubsection{4.6 Practical
consequences}\label{practical-consequences-2}}

Four framework-grounded predictions for perturbation prediction
methodology:

(P7) Transcriptome-only foundation models will not exceed the
conditional-mean baseline for perturbation prediction, regardless of
model scale or training data size. The ceiling is information-theoretic
and not addressable through more pretraining.

(P8) Methods incorporating chromatin accessibility data (paired
Perturb-seq + ATAC-seq) will systematically outperform RNA-only methods
on perturbation prediction. The improvement should be largest for
perturbations affecting chromatin-state transitions (knockdowns of
chromatin modifiers, transcription factors with autoregulatory loops).

(P9) Methods incorporating prior knowledge (knowledge graphs, motif
libraries, regulatory network priors) will systematically outperform
pure data-driven methods. The improvement scales with how well the prior
captures the actual hidden-sector structure relevant to the prediction
task.

(P10) Compositional generalization (predicting double perturbations from
single-perturbation training) will be systematically harder for
perturbations co-localized at the chromatin level than for perturbations
acting on disjoint pathways. This is testable through benchmarking on
Norman et al.'s combinatorial dataset stratified by chromatin
co-localization of the perturbed gene pairs.

The fourth prediction is particularly diagnostic. If it holds, it
provides direct empirical confirmation that the framework's structural
picture (chromatin as hidden sector mediating perturbation interactions)
is the right explanation for compositional generalization failure modes.

\begin{center}\rule{0.5\linewidth}{0.5pt}\end{center}

\hypertarget{multi-omics-integration}{%
\subsection{5. Multi-Omics Integration}\label{multi-omics-integration}}

\hypertarget{the-multi-omics-integration-paradigm}{%
\subsubsection{5.1 The multi-omics integration
paradigm}\label{the-multi-omics-integration-paradigm}}

Single-cell multi-omics technologies measure multiple molecular
modalities in the same cell: G\&T-seq (genome + transcriptome),
scNMT-seq (chromatin accessibility + transcriptome + methylation),
CITE-seq (transcriptome + surface protein), QuRIE-seq (transcriptome +
intracellular phosphoprotein), and increasingly multimodal Perturb-seq
variants. The data architecture is: \(N\) cells × \(M\) modalities ×
\(G_m\) features per modality, with the modalities measuring different
aspects of cellular state on different timescales.

The standard computational task is integration: combining information
across modalities to produce unified cell embeddings, joint
trajectories, and integrated regulatory networks. The principal methods
are Multi-Omics Factor Analysis (MOFA {[}36{]} and MOFA+ {[}37{]}),
Seurat's weighted nearest neighbor (WNN) {[}38{]}, LIGER {[}39{]}, and
totalVI {[}40{]}. These methods share a common assumption: different
modalities are parallel measurements of the same underlying cell state,
and integration combines them into a single low-dimensional latent
space.

\hypertarget{the-timescale-problem}{%
\subsubsection{5.2 The timescale problem}\label{the-timescale-problem}}

The parallel-measurements assumption fails when modalities measure
different memory layers operating on different timescales. The QuRIE-seq
paper {[}8{]} explicitly states: ``The timescale of (phospho)proteome
changes falls into a minute regime, whereas most significant
transcriptome changes are observed after hours.'' The same cells,
measured at the same instant, contain information from different points
in their dynamical history --- phosphoprotein measurements reflect
minute-scale signaling state, transcriptome measurements reflect
hour-scale gene expression state, and methylation measurements reflect
day-to-week-scale chromatin state.

The framework's perspective on this is direct: each modality measures a
different memory layer with a different \(\tau_B\). The visible-sector /
hidden-sector partition is modality-dependent. From the perspective of
phosphoprotein dynamics, the transcriptome state is part of the slow
hidden sector (chromatin sets transcriptional rates, which set protein
production rates). From the perspective of methylation dynamics, the
transcriptome is part of the fast visible sector (transcription happens
fast on the methylation timescale). The static factor-analysis
assumption --- that all modalities reflect the same cell state ---
misses this hierarchical structure.

The empirical consequences are documentable. Static integration methods
produce embeddings that are dominated by the modality with the largest
dynamic range during the measurement window --- typically transcriptome
--- and the other modalities contribute mostly as confirmation of the
transcriptome-derived structure rather than as independent information.
The integration improves accuracy modestly compared to single-modality
methods, but does not fully exploit the temporal information present in
the multi-modal data.

\hypertarget{emerging-timescale-aware-methods}{%
\subsubsection{5.3 Emerging timescale-aware
methods}\label{emerging-timescale-aware-methods}}

Recent methodological work has begun to address the timescale problem
explicitly. The QuRIE-seq paper {[}8{]} used time-series sampling
combined with MOFA+ factors to capture dynamics across modalities. Heyne
et al.~(2024) {[}41{]} introduced pseudotemporal integration of
single-cell RNA-Seq and single-cell proteomics by mass spectrometry,
building cell-state trajectories separately for each modality and
aligning them through shared markers. The approach yields
``time-dependent, nonlinear relationships between transcript and protein
levels'' that static methods cannot recover.

The chromatin velocity literature (§2) is a special case of
timescale-aware multi-omics integration: MultiVelo {[}17{]} explicitly
models chromatin and transcription as coupled dynamical systems with
different switch times. The success of MultiVelo over scVelo is the
timescale-aware version of the framework's prediction --- using temporal
hierarchy in the integration produces better dynamical estimates.

The framework predicts further methodological development along these
lines will succeed. Specifically, integration methods that include
explicit lag structure between modalities (transcriptome at time \(t\)
reflecting chromatin state at time \(t - \tau_B\)) should outperform
methods that treat modalities as instantaneous parallel measurements.
The lag structure should be biologically calibrated to the
characteristic \(\tau_B\) of each modality.

\hypertarget{the-hierarchy-of-biological-memory-layers}{%
\subsubsection{5.4 The hierarchy of biological memory
layers}\label{the-hierarchy-of-biological-memory-layers}}

The framework provides a specific prediction for the timescale hierarchy
across biological modalities, ordered from fastest to slowest:

\begin{longtable}[]{@{}
  >{\raggedright\arraybackslash}p{(\columnwidth - 4\tabcolsep) * \real{0.3333}}
  >{\raggedright\arraybackslash}p{(\columnwidth - 4\tabcolsep) * \real{0.3333}}
  >{\raggedright\arraybackslash}p{(\columnwidth - 4\tabcolsep) * \real{0.3333}}@{}}
\toprule\noalign{}
\begin{minipage}[b]{\linewidth}\raggedright
Modality
\end{minipage} & \begin{minipage}[b]{\linewidth}\raggedright
Typical \(\tau_B\)
\end{minipage} & \begin{minipage}[b]{\linewidth}\raggedright
Memory layer
\end{minipage} \\
\midrule\noalign{}
\endhead
\bottomrule\noalign{}
\endlastfoot
Phosphoproteins & seconds--minutes & Signaling state \\
Metabolites & seconds--minutes & Metabolic state \\
Transcriptome & minutes--hours & Gene expression state \\
Histone acetylation & minutes--hours & Active chromatin state \\
Histone methylation & hours--days & Stable chromatin state \\
Proteome & hours--days & Translation state \\
Chromatin accessibility & hours--days & Regulatory element state \\
DNA methylation & days--weeks & Stable epigenetic state \\
Genome (sequence) & \(\to \infty\) & Heritable state \\
Chromatin 3D architecture (Hi-C) & hours--generations & Genome
organization \\
\end{longtable}

These timescales are approximate and tissue-dependent, but the ordering
is robust. The framework predicts that integration methods that respect
this ordering --- treating slower modalities as constraints on faster
modalities --- will systematically outperform methods that treat them
symmetrically.

\hypertarget{static-factor-analysis-under-the-framework}{%
\subsubsection{5.5 Static factor analysis under the
framework}\label{static-factor-analysis-under-the-framework}}

MOFA and MOFA+ are widely used and produce useful cell embeddings. They
do not, however, capture the temporal hierarchy across modalities. The
framework's prediction is that MOFA-style methods extract the
\emph{static covariance structure} shared across modalities --- what is
common to all measurements of a cell at one instant --- but miss the
\emph{dynamical relationships} between modalities --- how chromatin
state at one time predicts transcriptome state at a later time. The
static covariance is informative but incomplete; the dynamical
relationships contain additional information that static methods cannot
access.

This is testable. A direct comparison of MOFA+ integration with a
timescale-aware method on time-series data should show the
timescale-aware method recovering dynamical relationships that MOFA
misses. The QuRIE-seq paper {[}8{]} and subsequent time-series
multi-omics studies provide the data substrate for such comparisons; the
comparison has not been systematically performed across the
methodological literature.

\hypertarget{practical-consequences-3}{%
\subsubsection{5.6 Practical
consequences}\label{practical-consequences-3}}

Three framework-grounded predictions for multi-omics integration:

(P11) Static factor-analysis integration methods (MOFA, MOFA+, LIGER,
totalVI) systematically underutilize the temporal information present in
multi-modal data. The performance gap relative to timescale-aware
methods should be largest in dynamic biological processes (perturbation
responses, differentiation) and smallest in steady-state systems.

(P12) Timescale-aware integration methods will systematically outperform
static methods on prediction tasks involving cellular dynamics. The
improvement should scale with the ratio between the slowest and fastest
modality timescales in the data: large ratios (transcriptome +
methylation) should show large improvements; small ratios (transcriptome
+ chromatin accessibility, since both operate on hour-scale timescales)
should show smaller improvements.

(P13) The optimal modality combination for a given biological question
depends on the characteristic timescale of the process. Studies of fast
signaling dynamics should prioritize phosphoproteome + transcriptome;
studies of stable cell-fate decisions should prioritize methylome +
transcriptome; studies of developmental trajectories should include
chromatin architecture (Hi-C) as the slowest layer.

The framework provides a structural rationale for the multi-omics
integration practices that bioinformaticians are converging on through
empirical experience. Rather than treating multi-omics as a generic
improvement, the framework predicts which modality combinations matter
for which biological questions, and which integration methods will
succeed.

\begin{center}\rule{0.5\linewidth}{0.5pt}\end{center}

\hypertarget{cancer-methylome-classification}{%
\subsection{6. Cancer Methylome
Classification}\label{cancer-methylome-classification}}

\hypertarget{the-methylation-based-tumor-classification-paradigm}{%
\subsubsection{6.1 The methylation-based tumor classification
paradigm}\label{the-methylation-based-tumor-classification-paradigm}}

DNA methylation profiling has emerged over the past decade as a central
diagnostic tool in neuro-oncology and increasingly across cancer types.
The Heidelberg CNS Tumor Classifier, introduced by Capper et al.~(Nature
2018) {[}42{]} with version v11 trained on 2,801 reference samples,
classified central nervous system tumors into 91 methylation classes.
Version 12.8 (Cancer Cell 2025) {[}43{]} expanded this to 184 subclasses
trained on 7,495 methylation profiles, achieving 95\% subclass-level
accuracy with well-calibrated probabilistic scores. Methylation-based
classification is now incorporated into the WHO Classification of CNS
Tumors (CNS5, 2021) as a desirable or essential method for accurate
diagnosis of multiple tumor types {[}44{]}.

The classifier reads DNA methylation patterns at hundreds of thousands
of CpG sites (Illumina EPIC 850K/930K array) and assigns the sample to a
reference tumor class based on similarity to a curated training set. The
methylation signature is informative because, as Capper et al.~(2018)
{[}42{]} explicitly note, ``the genome-wide DNA methylation pattern (the
methylome) likely represents a combination of both the cell of origin
and somatically acquired DNA methylation changes.'' The cell-of-origin
component dominates the classification: methylation patterns reliably
identify the developmental lineage from which the tumor arose, even when
the tumor's morphology, histology, or transcriptional state has diverged
substantially from the lineage's normal pattern.

\hypertarget{methylome-arrest-in-cancer}{%
\subsubsection{6.2 Methylome arrest in
cancer}\label{methylome-arrest-in-cancer}}

A critical observation in the methylation-classification literature is
that DNA methylation patterns are \emph{relatively stable} in cancer
cells once established. Capper et al.~(2018) {[}42{]}: ``the dynamics of
the methylome crucial for cellular differentiation in non-transformed
cells seem to be somewhat arrested in cancer cells and DNA methylation
signatures likely remain relatively stable during the course of
disease.'' The methylome is not dynamic in the same way as the
transcriptome or proteome; it is stable enough to function as a reliable
diagnostic marker across disease progression, treatment, and recurrence.

This stability is what makes methylation classification work clinically.
If cancer methylomes were as dynamic as cancer transcriptomes, a single
methylation profile could not be reliably associated with a tumor class
--- the profile would shift with disease progression and treatment, just
as transcriptional profiles do. The methylome's relative stability
across disease course is the practical foundation of the entire
methylation-classification paradigm.

\hypertarget{the-substratum-emergent-operator-distinction-in-cancer}{%
\subsubsection{6.3 The substratum-emergent operator distinction in
cancer}\label{the-substratum-emergent-operator-distinction-in-cancer}}

The framework's predicted pattern for cancer is the substratum-emergent
operator distinction {[}Medicine §9.4{]} applied to oncology: the
substratum-level memory operator (DNA methylation pattern) is preserved
with high fidelity reflecting the cell of origin, while the
emergent-level operator (cell behavior, proliferation, differentiation,
transcriptional phenotype) is dysregulated. The two operators dissociate
in cancer in a specific direction: substratum preserved, emergent
dysregulated.

This is the structural pattern the methylation-classification field has
empirically established. The 95\% subclass-level accuracy of
methylation-based classification reflects the high fidelity of the
substratum-level memory; the fact that the same methylation patterns
associate with diverse histological appearances, transcriptional states,
and clinical behaviors reflects the dysregulation of the emergent
operator. The framework's prediction matches the field's diagnostic
paradigm precisely.

This same pattern operates across scales: sphalerons and baryon-number
violation in particle physics, iPSC reprogramming residuals in stem cell
biology {[}Medicine §9.4{]}, and methylation-classifier-based diagnosis
in oncology are all instances of the broader substratum-emergent
operator distinction observed in foundational physics {[}SM §8.7{]} ---
where operators conserved at the substratum (lattice) level can be
violated at the emergent (effective field theory) level. The pattern is
unified across scales: the substratum carries information that the
emergent operator can violate, and methylation classification in
oncology operationalizes this for clinical diagnostics.

\hypertarget{unclassifiable-tumors-and-substratum-disruption}{%
\subsubsection{6.4 Unclassifiable tumors and substratum
disruption}\label{unclassifiable-tumors-and-substratum-disruption}}

A specific empirical pattern in methylation classification is
informative for the framework. Drexler et al.~(2024) {[}45{]} document
that approximately 10--15\% of CNS tumor cases submitted to the
Heidelberg classifier are ``unclassifiable'' --- they do not match any
reference methylation class with confidence above the diagnostic
threshold (calibrated score \(\geq 0.9\)). These unclassifiable cases
have systematically distinct features: lower tumor purity, younger
patient age, recurrent tumor tissue post-radiotherapy, and significantly
worse clinical outcomes including longer time to treatment initiation
and less favorable survival in IDH wildtype glioblastomas.

The framework's prediction for unclassifiable tumors is that they are
cases where the substratum-level memory operator has been disrupted, not
just the emergent operator. The disrupted substratum produces a
methylation pattern that no longer reflects a single cell of origin, and
the classifier cannot map it to a reference class. The empirical
features of unclassifiable cases match this prediction: lower tumor
purity (contributions from multiple cell types altering the methylation
pattern), post-radiotherapy recurrence (DNA damage altering methylation
maintenance), and younger age (less time for methylation to stabilize
into a clear pattern).

This is testable beyond clinical features. The framework predicts that
unclassifiable tumors should show specifically altered methylation
patterns at maintenance sites (DNMT1 binding sites, hemi-methylated
regions) and at chromatin boundary elements, reflecting disruption of
the methylation maintenance machinery rather than tumor-type-specific
methylation differences. The hypothesis is testable in the existing
Heidelberg dataset by stratifying classifiable versus unclassifiable
cases on machinery-site versus content-site methylation patterns.

\hypertarget{methylation-classification-beyond-cns-tumors}{%
\subsubsection{6.5 Methylation classification beyond CNS
tumors}\label{methylation-classification-beyond-cns-tumors}}

The methylation-classification paradigm has been extended to other
cancer types: sarcoma classifiers {[}46{]}, leukemia classifiers, and
pan-cancer methylation classifiers {[}47{]}. The performance pattern is
consistent: high accuracy for tumor type identification, with the
methylation signature carrying cell-of-origin information that survives
the tumor's emergent dysregulation.

The framework predicts that this paradigm will continue to expand across
cancer types, and that the boundary cases --- cancer types where
methylation classification has lower accuracy --- will be those where
the substratum-level memory is more readily disrupted in the cancer
state. Highly mutated cancers with extensive copy number aberrations
(which can disrupt chromatin maintenance machinery) should be less
amenable to methylation classification than cancers with stable genomes.
This is a specific testable prediction about cross-cancer methylation
classifier performance.

\hypertarget{practical-consequences-4}{%
\subsubsection{6.6 Practical
consequences}\label{practical-consequences-4}}

Three framework-grounded predictions for cancer methylome
bioinformatics:

(P14) The methylation-classification paradigm will continue to extend to
cancer types where the substratum-level memory is preserved during
tumorigenesis. Cancers with stable genomes (most pediatric CNS tumors,
well-differentiated sarcomas) will achieve high classifier accuracy;
cancers with extensive chromatin maintenance machinery mutations will
achieve lower accuracy.

(P15) Unclassifiable tumors are not a single category but split into two
distinct groups: (i) cancers from rare cell-of-origin lineages not
represented in the training data (the standard interpretation), and (ii)
cancers where the substratum-level memory has been disrupted (the
framework's prediction). The two groups should have distinguishable
methylation patterns at maintenance versus content sites.

(P16) Methylation-based classifiers can be improved by including
chromatin maintenance machinery integrity (mutations or copy number
alterations in DNMT1, DNMT3A/B, TET1/2/3, etc.) as a covariate. The
framework predicts that tumors with intact maintenance machinery should
be more reliably classifiable than tumors with disrupted machinery, at
given tumor purity.

The third prediction is testable through reanalysis of the existing
Heidelberg cohort with maintenance machinery integrity as a
stratification variable.

\begin{center}\rule{0.5\linewidth}{0.5pt}\end{center}

\hypertarget{variant-pathogenicity-for-readerwritereraser-genes}{%
\subsection{7. Variant Pathogenicity for Reader/Writer/Eraser
Genes}\label{variant-pathogenicity-for-readerwritereraser-genes}}

\hypertarget{variant-effect-prediction-methods}{%
\subsubsection{7.1 Variant effect prediction
methods}\label{variant-effect-prediction-methods}}

Predicting the functional impact of missense variants on protein-coding
genes is a core task in clinical genetics. Methods range from
sequence-conservation methods (SIFT {[}48{]}, PolyPhen-2 {[}49{]})
through structure-aware methods (CADD {[}50{]}, REVEL {[}51{]}) to
AI-based methods (AlphaMissense {[}52{]}). AlphaMissense, introduced by
Cheng et al.~(Science 2023) {[}52{]}, applies the AlphaFold2
architecture {[}53{]} to variant effect prediction and achieves
state-of-the-art performance across multiple benchmarks, providing
pathogenicity scores for every possible single-amino-acid substitution
in the human proteome {[}54{]}.

The methods operate on a common assumption: variant pathogenicity is
determined by the variant's impact on the protein it directly affects.
Conservation features capture evolutionary intolerance to variation at
the residue; structural features capture the variant's disruption of
folding or function; AI features capture the integration of structural
and conservation information at a higher level. The methods do not, in
general, account for the variant's \emph{functional role} in the broader
regulatory architecture --- whether the affected gene is a substrate
gene, a reader of epigenetic marks, a writer that establishes marks, or
an eraser that removes marks.

\hypertarget{the-frameworks-prediction}{%
\subsubsection{7.2 The framework's
prediction}\label{the-frameworks-prediction}}

The Medicine companion paper {[}Medicine §10.7{]} establishes the
reader-writer disorder class: genetic disorders caused by
loss-of-function mutations in readers, writers, or erasers of epigenetic
memory (Rett syndrome via MECP2, Fragile X via FMR1, Angelman via UBE3A,
Kabuki via KMT2D, ATR-X via ATRX). These disorders have substantially
different phenotypic features than disorders affecting substrate genes
directly: wider therapeutic windows, demonstrated reversibility in
animal models (Guy et al.~2007 {[}55{]} for Rett), and graded severity
correlated with residual machinery function.

The framework predicts that variant pathogenicity prediction should be
systematically different for variants in reader/writer/eraser machinery
versus variants in substrate genes. Specifically:

\begin{enumerate}
\def\labelenumi{(\roman{enumi})}
\item
  Variants in machinery genes should show \emph{graded} pathogenicity
  across mutations within the same gene, reflecting the gene's
  dose-dependent function in maintaining substratum-level memory across
  the genome. Different MECP2 mutations produce different Rett severity
  precisely because they confer different degrees of residual reading
  function {[}56{]}.
\item
  The relationship between variant pathogenicity score and clinical
  severity should differ in slope between machinery genes and substrate
  genes. For substrate genes, loss-of-function correlates directly with
  phenotype severity (a missing enzyme produces a deficient pathway).
  For machinery genes, loss-of-function correlates with the \emph{number
  of downstream targets affected}, which can be partially compensated by
  residual machinery function --- producing a flatter
  pathogenicity-severity slope.
\item
  The relative performance of variant effect predictors on machinery
  genes versus substrate genes should differ. Methods trained
  predominantly on substrate genes (where the conservation-pathogenicity
  relationship is well-calibrated) should systematically underperform on
  machinery genes (where the relationship is mediated by downstream
  effects on chromatin state).
\end{enumerate}

\hypertarget{the-proposed-empirical-test}{%
\subsubsection{7.3 The proposed empirical
test}\label{the-proposed-empirical-test}}

The framework's predictions can be tested through reanalysis of existing
AlphaMissense predictions stratified by gene functional role. The
protocol:

\begin{enumerate}
\def\labelenumi{\arabic{enumi}.}
\item
  Classify genes into three categories: substrate genes (encoding
  enzymes, transporters, structural proteins, ligands, receptors),
  machinery genes (encoding epigenetic readers, writers, erasers,
  chromatin remodelers, transcription factors), and architectural genes
  (encoding scaffolds and connecting structures). Standard ontologies
  (GO Molecular Function, ChromBase, Eukaryotic Linear Motif database)
  provide the classification.
\item
  For each gene category, examine the distribution of AlphaMissense
  pathogenicity scores conditional on ClinVar pathogenicity assignment.
\item
  Test prediction (iii): machinery gene scores should show poorer
  correlation with ClinVar assignments than substrate gene scores, at
  matched evolutionary conservation level.
\item
  Test prediction (i): within machinery genes, the variance of
  pathogenicity scores for clinically observed variants should be higher
  than within substrate genes, reflecting the graded relationship
  between machinery function and phenotype.
\item
  Test prediction (ii): correlate variant pathogenicity score with
  phenotype severity (where available from clinical literature) and
  compare slopes across gene categories.
\end{enumerate}

The test is computationally tractable using existing public data
(AlphaMissense predictions on Zenodo, ClinVar, gene ontology databases).
Negative results would refute the framework's prediction; positive
results would establish a new structural pattern in variant effect
prediction with practical clinical implications.

\hypertarget{implications-for-clinical-variant-interpretation}{%
\subsubsection{7.4 Implications for clinical variant
interpretation}\label{implications-for-clinical-variant-interpretation}}

If the framework's predictions are confirmed, variant interpretation in
reader-writer disorders should incorporate gene functional role as a
covariate. Current clinical variant interpretation (ACMG/AMP guidelines
{[}57{]}) treats all genes symmetrically: variants are scored on
evolutionary conservation, computational predictions, segregation data,
functional studies, and population frequency, without distinguishing the
gene's functional role in the regulatory architecture. The framework
predicts this symmetric treatment systematically misclassifies variants
in machinery genes --- both overcalling and undercalling depending on
the specific gene's network position.

A revised approach would (i) flag machinery-gene variants for additional
scrutiny, recognizing that pathogenicity for these genes is mediated by
downstream effects rather than direct functional loss; (ii) recommend
functional assays measuring residual machinery activity rather than
relying on prediction scores alone; and (iii) interpret machinery-gene
variants in the context of the patient's epigenetic state, since the
same variant may produce different phenotypes depending on what
substrate it acts on. These recommendations align with emerging clinical
practice for Rett-class disorders {[}58{]} but have not been generalized
to the full machinery-gene category.

\hypertarget{practical-consequences-5}{%
\subsubsection{7.5 Practical
consequences}\label{practical-consequences-5}}

Three framework-grounded predictions for variant effect prediction
methodology:

(P17) AlphaMissense and similar methods systematically miscalibrate
pathogenicity scores for reader/writer/eraser genes relative to
substrate genes. The miscalibration is testable through stratified
reanalysis of existing predictions.

(P18) Variant effect prediction methods trained specifically on
machinery-gene variant data should outperform general-purpose methods on
this gene category. The methodological recommendation is to develop
machinery-gene-specific predictors that incorporate the gene's network
position and downstream regulatory targets.

(P19) Clinical variant interpretation should incorporate gene functional
role (substrate vs.~machinery vs.~architectural) as a structural feature
affecting how computational pathogenicity scores are weighted. The
framework predicts machinery-gene variants require additional functional
assay evidence beyond computational scores.

The third prediction has direct clinical implications and is testable
through retrospective audit of diagnostic genetic testing outcomes
stratified by gene category.

\begin{center}\rule{0.5\linewidth}{0.5pt}\end{center}

\hypertarget{reader-writer-disorder-bioinformatics}{%
\subsection{8. Reader-Writer Disorder
Bioinformatics}\label{reader-writer-disorder-bioinformatics}}

\hypertarget{the-frameworks-reader-writer-disorder-class}{%
\subsubsection{8.1 The framework's reader-writer disorder
class}\label{the-frameworks-reader-writer-disorder-class}}

The Medicine paper {[}Medicine §10.7{]} introduces the reader-writer
disorder class: genetic disorders caused by loss-of-function mutations
in epigenetic readers (MECP2 in Rett syndrome, FMR1 in Fragile X),
writers (UBE3A in Angelman syndrome, KMT2D in Kabuki syndrome), or
remodelers (ATRX in ATR-X syndrome). These disorders share a structural
feature: the substratum-level memory operator (DNA methylation, histone
modifications) is largely preserved in patients, while the trace-out
machinery that translates substrate state into cellular function is
disrupted. The framework predicts that this structural feature produces
specific bioinformatic patterns that distinguish reader-writer disorders
from substrate-affecting genetic disorders.

\hypertarget{bioinformatic-patterns-in-rett-syndrome}{%
\subsubsection{8.2 Bioinformatic patterns in Rett
syndrome}\label{bioinformatic-patterns-in-rett-syndrome}}

Rett syndrome provides the most directly testable case. MECP2 is a
methyl-CpG-binding protein --- a reader of DNA methylation. Patients
carry loss-of-function MECP2 mutations; the DNA methylation substrate
they read is generally intact {[}59, 60{]}. The framework predicts: in
Rett patient samples, DNA methylation patterns should be approximately
normal (substratum preserved), while gene expression and chromatin state
should be altered (emergent operator disrupted).

The empirical literature is consistent with this prediction. Picard et
al.~(2021) {[}60{]} review the transcriptomic and epigenomic landscape
in Rett: while large-scale alterations of the epigenome have been
hypothesized, ``the impact of MeCP2 malfunction on global DNA
methylation and its large-scale effect in RTT has been poorly
investigated and is still elusive.'' Initial reports of compensatory DNA
methylation differences between Rett patients and controls in small
cohorts were \emph{not confirmed} in larger cohorts. The methylation
substrate appears largely preserved.

By contrast, transcriptional alterations in Rett are well-documented:
thousands of differentially expressed genes between MECP2-mutant and
wildtype neurons, with characteristic patterns affecting long genes
{[}61{]}, synaptic plasticity genes, and activity-dependent genes. The
emergent transcriptional operator is substantially disrupted while the
substrate is largely preserved --- exactly the framework's prediction.

\hypertarget{generalization-to-other-reader-writer-disorders}{%
\subsubsection{8.3 Generalization to other reader-writer
disorders}\label{generalization-to-other-reader-writer-disorders}}

The framework predicts that similar bioinformatic patterns hold across
the reader-writer disorder class:

\textbf{Fragile X syndrome (FMR1).} FMR1 encodes an RNA-binding protein
that regulates translation of hundreds of target mRNAs. Loss of FMR1
should leave the RNA target landscape intact (substrate preserved) but
disrupt translation regulation (emergent operator). The empirical
pattern matches: FMR1-null cells show extensive translational
dysregulation while transcriptome changes are more modest {[}62, 63{]}.
The ``substrate'' in this case is the RNA target repertoire and ribosome
state; FMR1 is the reader.

\textbf{Angelman syndrome (UBE3A).} UBE3A is an E3 ubiquitin ligase
functioning as a writer of degradation marks (ubiquitin tags) on target
proteins. The protein substrate landscape is preserved; what is
disrupted is the writing of degradation marks. The framework predicts
protein-level dysregulation (accumulation of UBE3A targets) with
relatively preserved transcriptome and methylome. Empirical
confirmation: proteomic studies in Angelman models show selective
accumulation of UBE3A substrates {[}64{]}, with transcriptomic changes
secondary to the proteomic dysregulation.

\textbf{Kabuki syndrome (KMT2D).} KMT2D is a histone H3K4
methyltransferase --- a writer of activating histone marks at enhancers.
Loss of KMT2D should leave the chromatin substrate (nucleosome
positions, DNA methylation) intact but disrupt the deposition of
H3K4me1/me2/me3 marks. The framework predicts altered H3K4 methylation
patterns at affected enhancers with relatively preserved DNA methylation
and broader chromatin features. Empirical pattern: ChIP-seq studies of
KMT2D mutants show specific loss of H3K4 marks at enhancers, with other
chromatin features less affected {[}65{]}.

\textbf{ATR-X syndrome (ATRX).} ATRX is a chromatin remodeler that
maintains heterochromatin at telomeres, centromeres, and tandem repeat
sequences. Loss of ATRX should disrupt heterochromatin maintenance at
these sites while leaving euchromatin and most of the methylome intact.
The framework predicts site-specific chromatin defects with broad
methylome preservation. Empirical pattern: ATRX-deficient cells show
specific defects at heterochromatin maintenance sites with preservation
of broader chromatin architecture {[}66{]}.

The common pattern across the class: the substratum that the disrupted
reader/writer/eraser acts on is preserved; the emergent operator at the
level affected by the machinery is disrupted; other levels are
relatively preserved.

\hypertarget{differential-diagnosis-through-multi-omic-bioinformatics}{%
\subsubsection{8.4 Differential diagnosis through multi-omic
bioinformatics}\label{differential-diagnosis-through-multi-omic-bioinformatics}}

The framework's predicted patterns provide a bioinformatic differential
diagnostic. A patient with a developmental disorder of unknown cause can
be assessed for reader-writer-disorder pattern through paired multi-omic
profiling:

\begin{enumerate}
\def\labelenumi{\arabic{enumi}.}
\item
  \textbf{DNA methylation} (whole-genome bisulfite sequencing or EPIC
  array). Preserved methylation argues against substrate-disruption
  disorder; altered methylation argues for it.
\item
  \textbf{Histone modifications} at affected loci (ChIP-seq for H3K4me3,
  H3K27me3, H3K27ac). Loss of specific marks argues for writer
  disruption (KMT2D, EZH2, etc.); preservation argues against.
\item
  \textbf{Chromatin accessibility} (ATAC-seq). Altered accessibility at
  specific elements with preserved bulk accessibility argues for
  reader/writer disruption; bulk accessibility changes argue for global
  disruption.
\item
  \textbf{Transcriptome} (scRNA-seq). Characteristic dysregulation
  patterns (long genes for Rett; ribosomal targets for Fragile X; UBE3A
  target accumulation for Angelman) provide signature for specific
  disorder.
\end{enumerate}

The framework predicts that multi-omic patterns can distinguish
reader-writer disorders from substrate-affecting disorders even before
the causal mutation is identified --- a useful diagnostic capability for
patients with unclassified developmental syndromes.

\hypertarget{therapeutic-implications-and-bioinformatic-monitoring}{%
\subsubsection{8.5 Therapeutic implications and bioinformatic
monitoring}\label{therapeutic-implications-and-bioinformatic-monitoring}}

The Medicine paper {[}Medicine §10.7{]} establishes that reader-writer
disorders are reversibility-prone because the substrate has been
preserved. Gene therapy (AAV-based MECP2 replacement, UBE3A
reactivation) and pharmacological replacement (trofinetide for Rett,
FDA-approved 2023) target this property. Bioinformatic monitoring of
therapy outcomes should follow the framework's structural picture:

\begin{enumerate}
\def\labelenumi{(\roman{enumi})}
\item
  Therapeutic restoration of reader/writer function should produce
  \emph{transcriptome normalization} if the substrate is intact, with
  kinetics determined by the substrate's stability and the cell's
  ability to read the restored signal. Methylation patterns should
  remain stable through treatment.
\item
  Patients who do not respond to reader/writer restoration may have
  suffered substrate degradation in addition to machinery loss --- a
  more advanced disease state. Methylation profiling at baseline could
  predict therapeutic response: patients with preserved methylation
  should be more responsive than patients with altered methylation.
\item
  Long-term monitoring through multi-omic profiling should track
  substrate fidelity (methylation stability) alongside emergent operator
  function (transcriptome, behavior), with the framework predicting that
  substrate integrity is the prognostic marker for sustained response.
\end{enumerate}

\hypertarget{practical-consequences-6}{%
\subsubsection{8.6 Practical
consequences}\label{practical-consequences-6}}

Three framework-grounded predictions for reader-writer disorder
bioinformatics:

(P20) Reader-writer disorders should show characteristic multi-omic
patterns: substrate preservation (methylation, sequence) with selective
emergent operator disruption (transcriptome, chromatin, proteome at the
level of the disrupted machinery). The pattern is testable across the
class through standardized multi-omic profiling.

(P21) Differential diagnosis between reader-writer disorders and
substrate-affecting genetic disorders is achievable through multi-omic
patterns even without sequencing the causal gene. This enables
diagnostic workup for patients with unclassified developmental
syndromes.

(P22) Therapeutic response to reader/writer restoration therapy is
predicted by baseline substrate integrity. Methylation profiling at
treatment initiation should predict response duration and degree.

The second and third predictions have direct clinical applications and
motivate prospective multi-omic studies in reader-writer disorder
cohorts.

\begin{center}\rule{0.5\linewidth}{0.5pt}\end{center}

\hypertarget{synthesis}{%
\subsection{9. Synthesis}\label{synthesis}}

\hypertarget{five-domains-one-structural-explanation}{%
\subsubsection{9.1 Five domains, one structural
explanation}\label{five-domains-one-structural-explanation}}

The preceding sections have developed framework predictions across five
active methodological domains in single-cell bioinformatics. The
findings are summarized in Table 1.

\begin{longtable}[]{@{}
  >{\raggedright\arraybackslash}p{(\columnwidth - 8\tabcolsep) * \real{0.2000}}
  >{\raggedright\arraybackslash}p{(\columnwidth - 8\tabcolsep) * \real{0.2000}}
  >{\raggedright\arraybackslash}p{(\columnwidth - 8\tabcolsep) * \real{0.2000}}
  >{\raggedright\arraybackslash}p{(\columnwidth - 8\tabcolsep) * \real{0.2000}}
  >{\raggedright\arraybackslash}p{(\columnwidth - 8\tabcolsep) * \real{0.2000}}@{}}
\toprule\noalign{}
\begin{minipage}[b]{\linewidth}\raggedright
Domain
\end{minipage} & \begin{minipage}[b]{\linewidth}\raggedright
Documented failure mode
\end{minipage} & \begin{minipage}[b]{\linewidth}\raggedright
Framework's structural explanation
\end{minipage} & \begin{minipage}[b]{\linewidth}\raggedright
Predicted improvement
\end{minipage} & \begin{minipage}[b]{\linewidth}\raggedright
Empirical confirmation
\end{minipage} \\
\midrule\noalign{}
\endhead
\bottomrule\noalign{}
\endlastfoot
Trajectory inference (§2) & scVelo trajectory reversal on complex
datasets {[}3{]} & Chromatin hidden sector violates Markovian assumption
& Multi-omic velocity & MultiVelo {[}17{]}, mmVelo {[}18{]}, MoFlow
{[}19{]} confirm \\
GRN inference (§3) & Causal methods fail to outperform correlation
{[}5{]} & Marginal distribution insufficient for causal recovery &
Multi-modal methods + prior knowledge & SCENIC+ {[}23{]}, LINGER
{[}24{]} confirm \\
Perturbation prediction (§4) & Deep learning underperforms linear
baselines {[}6{]} & Information-theoretic ceiling on transcriptome-only
data & Knowledge graphs, multi-modal data & GEARS {[}32{]}, PerturbNet
{[}35{]} partially confirm \\
Multi-omics integration (§5) & Static factor analysis misses temporal
hierarchy {[}8{]} & Modalities measure different memory layers &
Timescale-aware integration & QuRIE-seq {[}8{]}, pseudotemporal methods
{[}41{]} emerging \\
Cancer methylome (§6) & (No failure; methylation classifies 95\%
accurately {[}43{]}) & Substratum preserved, emergent dysregulated in
cancer & Substratum-preservation as diagnostic basis & Methylation
classification paradigm operationalizes \\
\end{longtable}

Each row instantiates the same underlying structural picture:
visible-sector data is insufficient for the task because cellular
dynamics depends on hidden-sector memory; methods that include
hidden-sector information explicitly succeed; methods that try to
extract hidden-sector information from visible-sector data alone are
bounded by an information-theoretic ceiling.

\hypertarget{the-information-theoretic-ceiling}{%
\subsubsection{9.2 The information-theoretic
ceiling}\label{the-information-theoretic-ceiling}}

The framework's most consequential bioinformatic prediction is the
existence of an information-theoretic ceiling on transcriptome-only
methods. This is a strong claim with specific empirical content: no
method using only scRNA-seq data --- regardless of model architecture,
training data size, or computational scale --- can recover certain
quantities that depend on hidden-sector variables.

The quantities subject to the ceiling include: the cell's causal
regulatory network structure (§3), the response to unseen perturbations
(§4), the temporal lag relationships between chromatin and transcription
(§5), and the substratum-level memory operator (§6). For each, the
underlying biology integrates over hidden-sector variables before the
data is collected; the integration is irreversible from the
visible-sector data alone.

The ceiling is empirically demonstrated in the foundation-model
literature {[}6, 7{]}. Multiple foundation models trained on tens of
millions of cells with billions of parameters fail to outperform simple
baselines on tasks subject to the ceiling. This is the empirical
signature of a structural limit, not a methodological gap to be closed
by better algorithms.

The practical implication is direct: methodological development that
stays within the transcriptome-only paradigm faces diminishing returns.
Continued scaling of foundation models will produce marginal
improvements in tasks where the ceiling has not yet been reached (cell
type annotation in well-characterized tissues, gene expression
imputation), but cannot produce qualitative improvements in tasks
subject to the ceiling (causal GRN inference, perturbation prediction
for unseen drugs). The research investment justification for continuing
to scale foundation models depends on the assumption that more data and
compute will eventually overcome the apparent limits; the framework
predicts this assumption is wrong for ceiling-bounded tasks.

\hypertarget{the-roadmap-forward}{%
\subsubsection{9.3 The roadmap forward}\label{the-roadmap-forward}}

The framework's positive prediction is that methodological progress on
ceiling-bounded tasks requires architectural changes that incorporate
hidden-sector information. Three pathways are consistent with the
empirical pattern:

\textbf{Pathway 1: Multi-modal data.} Paired measurements of
transcriptome with chromatin accessibility (10x Multiome, scNMT-seq),
histone modifications (CITE-seq variants, Paired-Tag), methylation
(scNMT-seq), or chromatin architecture (scHi-C) provide direct
measurement of hidden-sector variables. Methods using paired data have
empirically demonstrated improvements across all five domains
(MultiVelo, SCENIC+, multi-modal perturbation methods, timescale-aware
integration, multi-omic classifiers).

\textbf{Pathway 2: Prior biological knowledge.} Knowledge graphs (Gene
Ontology, KEGG pathways, protein-protein interaction networks), motif
libraries (JASPAR, HOCOMOCO), and curated regulatory databases
(Regulome, ENCODE TFBS) encode hidden-sector structure that has been
determined through other experimental approaches. Methods incorporating
this prior knowledge as structural constraints (GEARS, SCENIC, LINGER)
outperform purely data-driven methods.

\textbf{Pathway 3: Explicit memory state representation.} Models that
maintain explicit chromatin-state variables alongside transcriptome
variables --- separate from the latent representations of deep learning
models --- can in principle capture non-Markovian dynamics. This
approach is less developed but theoretically aligned with the
framework's prediction. Mechanistic models in the MultiVelo style
(coupled ODEs for chromatin and transcription) are early examples.

The three pathways are complementary, not competing. The empirical
pattern suggests that combining them --- multi-modal data analyzed with
knowledge-graph priors using methods that maintain explicit memory state
--- represents the most promising methodological direction.

\hypertarget{the-substratum-emergent-operator-distinction-as-the-unifying-structural-feature}{%
\subsubsection{9.4 The substratum-emergent operator distinction as the
unifying structural
feature}\label{the-substratum-emergent-operator-distinction-as-the-unifying-structural-feature}}

A more compact summary of the synthesis: cellular bioinformatics
operates on systems that exhibit the substratum-emergent operator
distinction {[}SM §8.7, Medicine §9.4{]}. The substratum operators ---
chromatin state, DNA methylation, protein abundance, cell architecture
--- carry slow memory with high information content. The emergent
operators --- transcriptome state, instantaneous cell behavior --- carry
fast information with limited memory. Methods that read the substratum
operators (methylation classifiers, multi-omic methods) work; methods
that try to infer everything from emergent operators alone (RNA
velocity, transcriptome-only foundation models) fail in characteristic
patterns.

This is the same structural pattern that organizes the SM's treatment of
strong CP, baryon number conservation, sphalerons, and 't Hooft anomaly
matching {[}SM §8.7{]}, and that organizes Rett syndrome paradigm and
clinical epigenetic therapeutics {[}Medicine §9.4{]}. The pattern is
unified across foundational physics, cellular biology, and computational
biology --- instances of the same structural feature of trace-out
dynamics in C1--C3 systems.

The bioinformatic instantiation of the pattern has the empirical
advantage of being directly measurable: the substratum is accessible
through methylation arrays, ChIP-seq, and ATAC-seq; the emergent
operators are accessible through transcriptomics; their dissociation is
documentable in patient samples and tissue measurements. This is the
strongest empirical confirmation of the framework's structural pattern
across all its applications: the bioinformatics literature has
independently developed methods that operationalize the
substratum-emergent distinction without explicit reference to the
underlying framework, and the methods that work are precisely the ones
the framework predicts will work.

\begin{center}\rule{0.5\linewidth}{0.5pt}\end{center}

\hypertarget{open-questions-and-limits}{%
\subsection{10. Open Questions and
Limits}\label{open-questions-and-limits}}

\hypertarget{what-the-framework-predicts-but-has-not-been-tested}{%
\subsubsection{10.1 What the framework predicts but has not been
tested}\label{what-the-framework-predicts-but-has-not-been-tested}}

Several framework predictions in this paper are testable but not yet
empirically verified:

\begin{enumerate}
\def\labelenumi{(\roman{enumi})}
\item
  \textbf{Reader/writer/eraser variant pathogenicity stratification}
  (§7). The prediction that AlphaMissense and similar methods
  systematically miscalibrate machinery-gene variants is testable
  through reanalysis of existing predictions. The result, positive or
  negative, would be informative.
\item
  \textbf{Compositional generalization and chromatin co-localization}
  (§4.5). The prediction that double-perturbation prediction difficulty
  correlates with the chromatin co-localization of the perturbed genes
  can be tested using existing Perturb-seq datasets stratified by
  chromatin features.
\item
  \textbf{Unclassifiable tumor substratum disruption} (§6.4). The
  prediction that unclassifiable CNS tumors have specifically altered
  methylation at maintenance sites versus content sites is testable in
  the existing Heidelberg dataset.
\item
  \textbf{Timescale-aware integration superiority} (§5.6). The
  prediction that timescale-aware multi-omics methods systematically
  outperform static methods on dynamic biological processes is testable
  through controlled benchmarking on time-series multi-omic data.
\end{enumerate}

Each is a discrete empirical test with clear positive and negative
outcomes. The framework's bioinformatic predictions are falsifiable in a
specific operational sense.

\hypertarget{what-the-framework-does-not-predict}{%
\subsubsection{10.2 What the framework does not
predict}\label{what-the-framework-does-not-predict}}

The framework's predictions are limited to information-theoretic
constraints and structural patterns. It does not predict:

\begin{enumerate}
\def\labelenumi{(\roman{enumi})}
\item
  \textbf{Specific algorithm performance.} The framework predicts that
  multi-modal methods will outperform RNA-only methods, but does not
  predict which specific multi-modal method will be best for a given
  task. Algorithm-level performance depends on implementation details,
  hyperparameter tuning, and dataset-specific features that are outside
  the framework's scope.
\item
  \textbf{Quantitative effect sizes.} The framework predicts that
  knowledge-graph methods improve over data-driven methods, but does not
  predict the magnitude of improvement for a specific dataset. Effect
  sizes depend on the relevance of the available prior knowledge to the
  specific task.
\item
  \textbf{Failure modes outside the ceiling-bounded tasks.} The
  framework's information-theoretic ceiling applies to tasks requiring
  hidden-sector information. Other failure modes (batch effects,
  technical noise, dropouts, low cell numbers) are real and important
  but outside the framework's prediction scope.
\item
  \textbf{Specific molecular mechanisms.} The framework predicts
  structural patterns at the level of memory layers and information
  flows, not specific molecular pathways. Mechanistic biology continues
  to require domain-specific theoretical and experimental work.
\end{enumerate}

\hypertarget{open-structural-questions}{%
\subsubsection{10.3 Open structural
questions}\label{open-structural-questions}}

Three open questions remain at the structural level:

\textbf{Q1. What is the empirical \(\tau_B\) hierarchy across biological
systems?} The framework predicts that different memory layers have
different characteristic timescales, with consequences for which
modalities matter for which questions (§5.4). The actual \(\tau_B\)
values are tissue- and process-dependent, and have not been
systematically characterized. A program of measurements characterizing
\(\tau_B\) across cell types, tissues, and biological processes would
substantially refine the framework's bioinformatic predictions.

\textbf{Q2. How do different machinery genes contribute to the
trace-out?} The framework treats epigenetic readers, writers, and
erasers as components of the trace-out machinery, but the specific
contribution of each to the substratum-to-emergent translation is not
characterized in detail. Different machinery genes may produce different
patterns of substratum-emergent dissociation, with consequences for the
differential bioinformatic signatures of different reader-writer
disorders (§8.3).

\textbf{Q3. What is the relationship between the framework's
bioinformatic predictions and explicit kinetic models?} The framework's
predictions are information-theoretic and structural. Specific kinetic
models (mass-action models of gene regulation, stochastic gene
expression models, polymer models of chromatin) are mechanistic. The
relationship between the framework's structural picture and the
mechanistic models has not been worked out in detail. A formal
connection would clarify which aspects of mechanistic modeling
correspond to which framework structural features.

\hypertarget{limits}{%
\subsubsection{10.4 Limits}\label{limits}}

The framework's bioinformatic predictions have inherent limits:

\begin{enumerate}
\def\labelenumi{(\roman{enumi})}
\item
  \textbf{Scope limited to cellular systems satisfying C1--C3.} Systems
  where C1--C3 fail (rapidly equilibrating systems, low-capacity hidden
  sectors, decoupled visible-hidden dynamics) do not exhibit the
  framework's predicted patterns. Not every biological system has
  substantial \(\tau_B\) at relevant timescales; the framework's
  predictions apply where C1--C3 hold.
\item
  \textbf{Information-theoretic ceilings are upper bounds, not
  necessarily attainable.} A method may underperform the ceiling for
  many reasons unrelated to the framework: poor implementation,
  insufficient training data, wrong loss function. The framework
  predicts a ceiling but does not guarantee that any method will reach
  it.
\item
  \textbf{Empirical validation requires controlled benchmarking.} The
  framework's predictions are about structural patterns in method
  performance, which requires careful experimental design to test.
  Confounders (batch effects, dataset selection, evaluation metric
  choice) can obscure the patterns the framework predicts.
\item
  \textbf{The framework does not replace mechanistic biology.}
  Bioinformatic methods aim to extract biological insight from data, and
  biological insight requires understanding mechanisms. The framework
  provides a structural perspective on what mechanisms determine method
  performance, but does not substitute for the mechanistic understanding
  itself.
\end{enumerate}

These limits define the framework's scope of application in
computational biology. Within the scope --- systems with substantial
hidden-sector memory analyzed by methods that may or may not include
hidden-sector information --- the framework's predictions have empirical
content that has been confirmed across multiple methodological domains.

\begin{center}\rule{0.5\linewidth}{0.5pt}\end{center}

\hypertarget{acknowledgements}{%
\subsection{Acknowledgements}\label{acknowledgements}}

During the preparation of this work, the author used Claude Opus 4.7
(Anthropic) and Gemini 3.1 Pro (Google) to assist in drafting, refining
argumentation, and surveying the computational biology literature. The
author reviewed and edited all content and takes full responsibility for
the publication.

\begin{center}\rule{0.5\linewidth}{0.5pt}\end{center}

\hypertarget{references}{%
\subsection{References}\label{references}}

{[}1{]} A. Maybaum, ``The Incompleteness of Observation,'' (2026).

{[}2{]} A. Maybaum, ``Non-Markovian Dynamics in Biology and Medicine,''
companion paper (2026).

{[}3{]} G. Gorin, L. Pachter, N. Yosef, ``RNA velocity unraveled,''
\emph{PLOS Computational Biology} 18, e1010492 (2022).

{[}4{]} X. Yuan et al., ``Inferring gene regulatory networks from
single-cell multiome data using atlas-scale external data,''
\emph{Nature Biotechnology} (2024). doi:10.1038/s41587-024-02182-7.

{[}5{]} M. Passemiers et al., ``When Does Gene Regulatory Network
Inference Break? A Controlled Diagnostic Study of Causal and
Correlational Methods on Single-Cell Data,'' arXiv:2605.04930 (2025).

{[}6{]} C. Ahlmann-Eltze et al., ``Deep-learning-based gene perturbation
effect prediction does not yet outperform simple linear baselines,''
\emph{Nature Methods} (2025). doi:10.1038/s41592-025-02772-6.

{[}7{]} K. Kedzierska et al., ``Zero-shot evaluation reveals limitations
of single-cell foundation models,'' \emph{Genome Biology} 26, 101
(2025). doi:10.1186/s13059-025-03574-x.

{[}8{]} M. Reimegård et al., ``Single-cell intracellular epitope and
transcript detection reveals signal transduction dynamics'' (QuRIE-seq),
\emph{Cell Reports Methods} 1, 100165 (2021).

{[}9{]} G. La Manno et al., ``RNA velocity of single cells,''
\emph{Nature} 560, 494--498 (2018).

{[}10{]} V. Bergen et al., ``Generalizing RNA velocity to transient cell
states through dynamical modeling'' (scVelo), \emph{Nature
Biotechnology} 38, 1408--1414 (2020).

{[}11{]} M. Barile et al., ``Coordinated changes in gene expression
kinetics underlie both mouse and human erythroid maturation,''
\emph{Genome Biology} 22, 197 (2021).

{[}12{]} H. Cui et al., ``DeepVelo: Deep Learning extends RNA velocity
to multi-lineage systems with cell-specific kinetics,'' bioRxiv (2022).

{[}13{]} L. Liang et al., ``LatentVelo: structured latent gene
expression dynamics for RNA velocity estimation,'' bioRxiv (2022).

{[}14{]} A. Gayoso et al., ``Deep generative modeling of transcriptional
dynamics for RNA velocity analysis in single cells'' (veloVI),
\emph{Nature Methods} 21, 50--59 (2024).

{[}15{]} L. Liang et al., ``LatentVelo: structured latent gene
expression dynamics,'' bioRxiv (2022).

{[}16{]} A. Gayoso et al., ``veloVI: deep generative modeling of
transcriptional dynamics,'' \emph{Nature Methods} 21, 50--59 (2024).

{[}17{]} C. Li et al., ``Multi-omic single-cell velocity models
epigenome-transcriptome interactions and improves cell fate prediction''
(MultiVelo), \emph{Nature Biotechnology} 40, 1690--1698 (2022).

{[}18{]} N. Yamanaka et al., ``mmVelo: A deep generative model for
estimating cell state-dependent dynamics across multiple modalities,''
bioRxiv (2024). doi:10.1101/2024.12.11.628059.

{[}19{]} T. Tang et al., ``MoFlow: Multi-omic relay velocity modeling
uncovers dynamic chromatin-transcription regulation across cell
states,'' \emph{Nature Communications} (2025).
doi:10.1038/s41467-025-67259-6.

{[}20{]} X. Yuan et al., ``Inferring gene regulatory networks from
single-cell multiome data using atlas-scale external data,''
\emph{Nature Biotechnology} 42, 247--257 (2024).

{[}21{]} S. McCalla et al., ``Comprehensive Benchmarking of Single-Cell
Gene Regulatory Network Inference,'' \emph{Genome Biology} 24, 21
(2023).

{[}22{]} H. Nguyen et al., ``A comprehensive survey of regulatory
network inference methods using single cell RNA sequencing data,''
\emph{Briefings in Bioinformatics} 22, bbaa190 (2021).

{[}23{]} C. Bravo González-Blas et al., ``SCENIC+: single-cell multiomic
inference of enhancers and gene regulatory networks,'' \emph{Nature
Methods} 20, 1355--1367 (2023).

{[}24{]} X. Yuan et al., ``LINGER: Lifelong neural network for gene
regulation from single-cell multiome data,'' \emph{Nature Biotechnology}
42, 247--257 (2024).

{[}25{]} H. Cui et al., ``scGPT: toward building a foundation model for
single-cell multi-omics using generative AI,'' \emph{Nature Methods} 21,
1470--1480 (2024).

{[}26{]} C. V. Theodoris et al., ``Transfer learning enables predictions
in network biology'' (Geneformer), \emph{Nature} 618, 616--624 (2023).

{[}27{]} J. Doe et al., ``Systematic Evaluation of Single-Cell
Foundation Model Interpretability Reveals Attention Captures
Co-Expression Rather Than Unique Regulatory Signal,'' arXiv:2602.17532
(2026).

{[}28{]} J. M. Replogle et al., ``Mapping information-rich
genotype-phenotype landscapes with genome-scale Perturb-seq,''
\emph{Cell} 185, 2559--2575 (2022).

{[}29{]} T. M. Norman et al., ``Exploring genetic interaction manifolds
constructed from rich single-cell phenotypes,'' \emph{Science} 365,
786--793 (2019).

{[}30{]} M. Lotfollahi et al., ``Predicting cellular responses to
complex perturbations in high-throughput screens'' (CPA),
\emph{Molecular Systems Biology} 19, e11517 (2023).

{[}31{]} L. Hetzel et al., ``chemCPA: Predicting Cellular Responses to
Novel Drug Perturbations at a Single-Cell Resolution,'' NeurIPS (2022).

{[}32{]} Y. Roohani, K. Huang, J. Leskovec, ``Predicting transcriptional
outcomes of novel multigene perturbations with GEARS,'' \emph{Nature
Biotechnology} 42, 927--935 (2024).

{[}33{]} C. Bunne et al., ``Learning single-cell perturbation responses
using neural optimal transport,'' \emph{Nature Methods} 20, 1759--1768
(2023).

{[}34{]} M. Hao et al., ``Large-scale foundation model on single-cell
transcriptomics'' (scFoundation), \emph{Nature Methods} 21, 1481--1491
(2024).

{[}35{]} H. Yu et al., ``PerturbNet predicts single-cell responses to
unseen chemical and genetic perturbations,'' \emph{Molecular Systems
Biology} (2025).

{[}36{]} R. Argelaguet et al., ``Multi-Omics Factor Analysis (MOFA) ---
a framework for unsupervised integration of multi-omics data sets,''
\emph{Molecular Systems Biology} 14, e8124 (2018).

{[}37{]} R. Argelaguet et al., ``MOFA+: a statistical framework for
comprehensive integration of multi-modal single-cell data,''
\emph{Genome Biology} 21, 111 (2020).

{[}38{]} Y. Hao et al., ``Integrated analysis of multimodal single-cell
data'' (Seurat v4), \emph{Cell} 184, 3573--3587 (2021).

{[}39{]} J. D. Welch et al., ``Single-cell multi-omic integration
compares and contrasts features of brain cell identity'' (LIGER),
\emph{Cell} 177, 1873--1887 (2019).

{[}40{]} A. Gayoso et al., ``Joint probabilistic modeling of single-cell
multi-omic data with totalVI,'' \emph{Nature Methods} 18, 272--282
(2021).

{[}41{]} R. Heyne et al., ``Resolving Single-Cell Gene Expression by
Pseudotemporal Integration of Transcriptomic and Proteomic Datasets,''
\emph{PLOS Computational Biology} (2024).

{[}42{]} D. Capper et al., ``DNA methylation-based classification of
central nervous system tumours,'' \emph{Nature} 555, 469--474 (2018).

{[}43{]} D. Capper et al., ``Advancing CNS tumor diagnostics with
expanded DNA methylation-based classification'' (Heidelberg classifier
v12.8), \emph{Cancer Cell} (2025). doi:10.1016/j.ccell.2025.10.005.

{[}44{]} D. N. Louis et al., ``The 2021 WHO Classification of Tumors of
the Central Nervous System: a summary,'' \emph{Neuro-Oncology} 23,
1231--1251 (2021).

{[}45{]} R. Drexler et al., ``Unclassifiable CNS tumors in DNA
methylation-based classification: clinical challenges and prognostic
impact,'' \emph{Acta Neuropathologica Communications} 12, 9 (2024).

{[}46{]} C. Koelsche et al., ``Sarcoma classification by DNA methylation
profiling,'' \emph{Nature Communications} 12, 498 (2021).

{[}47{]} M. Renaud et al., ``Pan-cancer DNA methylation
classification,'' \emph{Genome Medicine} 14, 80 (2022).

{[}48{]} P. C. Ng, S. Henikoff, ``SIFT: predicting amino acid changes
that affect protein function,'' \emph{Nucleic Acids Research} 31,
3812--3814 (2003).

{[}49{]} I. A. Adzhubei et al., ``A method and server for predicting
damaging missense mutations'' (PolyPhen-2), \emph{Nature Methods} 7,
248--249 (2010).

{[}50{]} M. Kircher et al., ``A general framework for estimating the
relative pathogenicity of human genetic variants'' (CADD), \emph{Nature
Genetics} 46, 310--315 (2014).

{[}51{]} N. M. Ioannidis et al., ``REVEL: An Ensemble Method for
Predicting the Pathogenicity of Rare Missense Variants,'' \emph{American
Journal of Human Genetics} 99, 877--885 (2016).

{[}52{]} J. Cheng et al., ``Accurate proteome-wide missense variant
effect prediction with AlphaMissense,'' \emph{Science} 381, eadg7492
(2023).

{[}53{]} J. Jumper et al., ``Highly accurate protein structure
prediction with AlphaFold,'' \emph{Nature} 596, 583--589 (2021).

{[}54{]} L. Geistlinger et al., ``AlphaMissenseR: an integrated
framework for investigating missense mutations in human protein-coding
genes,'' \emph{Bioinformatics} (2025).
doi:10.1093/bioinformatics/btaf093.

{[}55{]} J. Guy et al., ``Reversal of neurological defects in a mouse
model of Rett syndrome,'' \emph{Science} 315, 1143--1147 (2007).

{[}56{]} D. C. Tarquinio et al., ``The course of awake breathing
disturbances across the lifespan in Rett syndrome,'' \emph{Brain \&
Development} 40, 515--529 (2018).

{[}57{]} S. Richards et al., ``Standards and guidelines for the
interpretation of sequence variants: a joint consensus recommendation of
the American College of Medical Genetics and Genomics and the
Association for Molecular Pathology,'' \emph{Genetics in Medicine} 17,
405--423 (2015).

{[}58{]} J. Kaufmann et al., ``Expert Recommendations for Trofinetide
Use in Individuals with Rett Syndrome,'' \emph{Pediatric Neurology}
(2024).

{[}59{]} J. Qian et al., ``Multiplex epigenome editing of MECP2 to
rescue Rett syndrome neurons,'' \emph{Science Translational Medicine}
15, eadd4666 (2023).

{[}60{]} N. Picard, M. Fagiolini, ``Transcriptomic and Epigenomic
Landscape in Rett Syndrome,'' \emph{Genes} 12, 967 (2021).

{[}61{]} H. W. Gabel et al., ``Disruption of DNA-methylation-dependent
long gene repression in Rett syndrome,'' \emph{Nature} 522, 89--93
(2015).

{[}62{]} J. C. Darnell et al., ``FMRP stalls ribosomal translocation on
mRNAs linked to synaptic function and autism,'' \emph{Cell} 146,
247--261 (2011).

{[}63{]} M. R. Sawicka et al., ``FMRP has a cell-type-specific role in
CA1 pyramidal neurons to regulate autism-related transcripts and
circadian memory,'' \emph{eLife} 8, e46919 (2019).

{[}64{]} Y. Yi et al., ``An autism-linked mutation disables
phosphorylation control of UBE3A,'' \emph{Cell} 162, 795--807 (2015).

{[}65{]} R. C. Cuellar et al., ``KMT2D and chromatin in development and
disease,'' \emph{Genes \& Development} 36, 1083--1100 (2022).

{[}66{]} G. Goldberg et al., ``Distinct factors control histone variant
H3.3 localization at specific genomic regions,'' \emph{Cell} 140,
678--691 (2010).

\hypertarget{the-frameworks-reader-writer-disorder-class-1}{%
\subsubsection{8.1 The framework's reader-writer disorder
class}\label{the-frameworks-reader-writer-disorder-class-1}}

The Medicine paper {[}Medicine §10.7{]} introduces the reader-writer
disorder class: genetic disorders caused by loss-of-function mutations
in epigenetic readers (MECP2 in Rett syndrome, FMR1 in Fragile X),
writers (UBE3A in Angelman syndrome, KMT2D in Kabuki syndrome), or
remodelers (ATRX in ATR-X syndrome). These disorders share a structural
feature: the substratum-level memory operator (DNA methylation, histone
modifications) is largely preserved in patients, while the trace-out
machinery that translates substrate state into cellular function is
disrupted. The framework predicts that this structural feature produces
specific bioinformatic patterns that distinguish reader-writer disorders
from substrate-affecting genetic disorders.

\hypertarget{bioinformatic-patterns-in-rett-syndrome-1}{%
\subsubsection{8.2 Bioinformatic patterns in Rett
syndrome}\label{bioinformatic-patterns-in-rett-syndrome-1}}

Rett syndrome provides the most directly testable case. MECP2 is a
methyl-CpG-binding protein --- a reader of DNA methylation. Patients
carry loss-of-function MECP2 mutations; the DNA methylation substrate
they read is generally intact {[}59, 60{]}. The framework predicts: in
Rett patient samples, DNA methylation patterns should be approximately
normal (substratum preserved), while gene expression and chromatin state
should be altered (emergent operator disrupted).

The empirical literature is consistent with this prediction. Picard et
al.~(2021) {[}60{]} review the transcriptomic and epigenomic landscape
in Rett: while large-scale alterations of the epigenome have been
hypothesized, ``the impact of MeCP2 malfunction on global DNA
methylation and its large-scale effect in RTT has been poorly
investigated and is still elusive.'' Initial reports of compensatory DNA
methylation differences between Rett patients and controls in small
cohorts were \emph{not confirmed} in larger cohorts. The methylation
substrate appears largely preserved.

By contrast, transcriptional alterations in Rett are well-documented:
thousands of differentially expressed genes between MECP2-mutant and
wildtype neurons, with characteristic patterns affecting long genes
{[}61{]}, synaptic plasticity genes, and activity-dependent genes. The
emergent transcriptional operator is substantially disrupted while the
substrate is largely preserved --- exactly the framework's prediction.

\hypertarget{generalization-to-other-reader-writer-disorders-1}{%
\subsubsection{8.3 Generalization to other reader-writer
disorders}\label{generalization-to-other-reader-writer-disorders-1}}

The framework predicts that similar bioinformatic patterns hold across
the reader-writer disorder class:

\textbf{Fragile X syndrome (FMR1).} FMR1 encodes an RNA-binding protein
that regulates translation of hundreds of target mRNAs. Loss of FMR1
should leave the RNA target landscape intact (substrate preserved) but
disrupt translation regulation (emergent operator). The empirical
pattern matches: FMR1-null cells show extensive translational
dysregulation while transcriptome changes are more modest {[}62, 63{]}.
The ``substrate'' in this case is the RNA target repertoire and ribosome
state; FMR1 is the reader.

\textbf{Angelman syndrome (UBE3A).} UBE3A is an E3 ubiquitin ligase
functioning as a writer of degradation marks (ubiquitin tags) on target
proteins. The protein substrate landscape is preserved; what is
disrupted is the writing of degradation marks. The framework predicts
protein-level dysregulation (accumulation of UBE3A targets) with
relatively preserved transcriptome and methylome. Empirical
confirmation: proteomic studies in Angelman models show selective
accumulation of UBE3A substrates {[}64{]}, with transcriptomic changes
secondary to the proteomic dysregulation.

\textbf{Kabuki syndrome (KMT2D).} KMT2D is a histone H3K4
methyltransferase --- a writer of activating histone marks at enhancers.
Loss of KMT2D should leave the chromatin substrate (nucleosome
positions, DNA methylation) intact but disrupt the deposition of
H3K4me1/me2/me3 marks. The framework predicts altered H3K4 methylation
patterns at affected enhancers with relatively preserved DNA methylation
and broader chromatin features. Empirical pattern: ChIP-seq studies of
KMT2D mutants show specific loss of H3K4 marks at enhancers, with other
chromatin features less affected {[}65{]}.

\textbf{ATR-X syndrome (ATRX).} ATRX is a chromatin remodeler that
maintains heterochromatin at telomeres, centromeres, and tandem repeat
sequences. Loss of ATRX should disrupt heterochromatin maintenance at
these sites while leaving euchromatin and most of the methylome intact.
The framework predicts site-specific chromatin defects with broad
methylome preservation. Empirical pattern: ATRX-deficient cells show
specific defects at heterochromatin maintenance sites with preservation
of broader chromatin architecture {[}66{]}.

The common pattern across the class: the substratum that the disrupted
reader/writer/eraser acts on is preserved; the emergent operator at the
level affected by the machinery is disrupted; other levels are
relatively preserved.

\hypertarget{differential-diagnosis-through-multi-omic-bioinformatics-1}{%
\subsubsection{8.4 Differential diagnosis through multi-omic
bioinformatics}\label{differential-diagnosis-through-multi-omic-bioinformatics-1}}

The framework's predicted patterns provide a bioinformatic differential
diagnostic. A patient with a developmental disorder of unknown cause can
be assessed for reader-writer-disorder pattern through paired multi-omic
profiling:

\begin{enumerate}
\def\labelenumi{\arabic{enumi}.}
\item
  \textbf{DNA methylation} (whole-genome bisulfite sequencing or EPIC
  array). Preserved methylation argues against substrate-disruption
  disorder; altered methylation argues for it.
\item
  \textbf{Histone modifications} at affected loci (ChIP-seq for H3K4me3,
  H3K27me3, H3K27ac). Loss of specific marks argues for writer
  disruption (KMT2D, EZH2, etc.); preservation argues against.
\item
  \textbf{Chromatin accessibility} (ATAC-seq). Altered accessibility at
  specific elements with preserved bulk accessibility argues for
  reader/writer disruption; bulk accessibility changes argue for global
  disruption.
\item
  \textbf{Transcriptome} (scRNA-seq). Characteristic dysregulation
  patterns (long genes for Rett; ribosomal targets for Fragile X; UBE3A
  target accumulation for Angelman) provide signature for specific
  disorder.
\end{enumerate}

The framework predicts that multi-omic patterns can distinguish
reader-writer disorders from substrate-affecting disorders even before
the causal mutation is identified --- a useful diagnostic capability for
patients with unclassified developmental syndromes.

\hypertarget{therapeutic-implications-and-bioinformatic-monitoring-1}{%
\subsubsection{8.5 Therapeutic implications and bioinformatic
monitoring}\label{therapeutic-implications-and-bioinformatic-monitoring-1}}

The Medicine paper {[}Medicine §10.7{]} establishes that reader-writer
disorders are reversibility-prone because the substrate has been
preserved. Gene therapy (AAV-based MECP2 replacement, UBE3A
reactivation) and pharmacological replacement (trofinetide for Rett,
FDA-approved 2023) target this property. Bioinformatic monitoring of
therapy outcomes should follow the framework's structural picture:

\begin{enumerate}
\def\labelenumi{(\roman{enumi})}
\item
  Therapeutic restoration of reader/writer function should produce
  \emph{transcriptome normalization} if the substrate is intact, with
  kinetics determined by the substrate's stability and the cell's
  ability to read the restored signal. Methylation patterns should
  remain stable through treatment.
\item
  Patients who do not respond to reader/writer restoration may have
  suffered substrate degradation in addition to machinery loss --- a
  more advanced disease state. Methylation profiling at baseline could
  predict therapeutic response: patients with preserved methylation
  should be more responsive than patients with altered methylation.
\item
  Long-term monitoring through multi-omic profiling should track
  substrate fidelity (methylation stability) alongside emergent operator
  function (transcriptome, behavior), with the framework predicting that
  substrate integrity is the prognostic marker for sustained response.
\end{enumerate}

\hypertarget{practical-consequences-7}{%
\subsubsection{8.6 Practical
consequences}\label{practical-consequences-7}}

Three framework-grounded predictions for reader-writer disorder
bioinformatics:

(P20) Reader-writer disorders should show characteristic multi-omic
patterns: substrate preservation (methylation, sequence) with selective
emergent operator disruption (transcriptome, chromatin, proteome at the
level of the disrupted machinery). The pattern is testable across the
class through standardized multi-omic profiling.

(P21) Differential diagnosis between reader-writer disorders and
substrate-affecting genetic disorders is achievable through multi-omic
patterns even without sequencing the causal gene. This enables
diagnostic workup for patients with unclassified developmental
syndromes.

(P22) Therapeutic response to reader/writer restoration therapy is
predicted by baseline substrate integrity. Methylation profiling at
treatment initiation should predict response duration and degree.

The second and third predictions have direct clinical applications and
motivate prospective multi-omic studies in reader-writer disorder
cohorts.

\begin{center}\rule{0.5\linewidth}{0.5pt}\end{center}

\end{document}
